%!TEX program = xelatex
% 使用 ctexart 文类，UTF-8 编码
\documentclass[UTF8,11pt, a4paper, oneside, fontset=none]{ctexbook}
\usepackage{ctex}
\setCJKmainfont{Songti SC}
\setCJKsansfont{Hiragino Sans GB}
\usepackage{caption}
\usepackage{subfigure}
\usepackage{setspace}
\usepackage{tipa}
\usepackage{longtable}
\usepackage{amsmath, amsthm, amssymb, bm, graphicx, hyperref, mathrsfs, geometry}



\title{数学分析习题课作业\\2023--2024学年按周整理}
\author{吉浩洋}
\date{}
\linespread{1.1}% 行间距
\newtheorem{theorem}{定理}[chapter]
\newtheorem{definition}[theorem]{定义}
\newtheorem{lemma}[theorem]{引理}
\newtheorem{corollary}[theorem]{推论}
\newtheorem{example}[theorem]{例}
\newtheorem{proposition}[theorem]{命题}
\newtheorem{formula}[theorem]{公式}
\newtheorem{remark}[theorem]{注记}
\newtheorem{question}[theorem]{问题}


\geometry{left=2cm, right= 2cm, bottom=3cm} %页面宽度

\allowdisplaybreaks



\begin{document}

\maketitle

\tableofcontents

\clearpage

\chapter{秋季学期}

\clearpage

\section{第一周}

\subsection*{题目}

\begin{itemize}
\item[1.] 第一节 1(2), 2(3), 3, 第二节 4.
\item[2] 利用确界原理证明以下阿基米德原理及其推论:\\
\noindent
{\it 阿基米德原理:} 设$h>0$为任一实数, 对每个$x \in \mathbb R$, 存在唯一的$k \in \mathbb Z$使得$(k-1)h \leq x < kh$. 

\noindent
{\it 阿基米德原理的推论:}
\begin{itemize}
\item[(1)] 对任意的$\epsilon >0$, 存在$n \in \mathbb N$使得$0 < 1/n < \epsilon$.
\item[(2)] 设$x \geq 0$, 若对所有$n \in \mathbb N$都有$x < 1/n$, 则$x =0$.
\item[(3)] 对于任意的$a, b \in \mathbb R$, 存在有理数$r$使得$a < r < b$. (有理数的稠密性)
\item[(4)] 对于任意的$x \in \mathbb R$, 存在唯一的整数$k$使得$k \leq x <k+1$. (另一形式的阿基米德原理)
\end{itemize}

\end{itemize}

\clearpage

\subsection*{原稿答案}

\noindent\emph{以下为按题干归入本周的原稿部分参考答案。}

\begin{itemize}

\item[(1)] 反证. 假设$ax$是有理数, 因为$a$是不为0的有理数, 故$\frac{ax}{a} = x$是有理数, 矛盾.

\item[(2)] 原不等式的解首先必须满足以下不等式组
\[
\begin{cases}
x-1 \geq 0,\\
2x-1 \geq 0,\\
3x-2 \geq 0.
\end{cases}
\]
解得$x \geq 1$. 而当$x \geq 1$时, 原不等式左端$\sqrt{x-1} - \sqrt{2x-1} <0$, 故无解.

\item[(3)] 反证. 假设$a \neq b$, 则$a -b \neq 0$. 设$\eta = |a-b|>0$. 取$\epsilon = \frac{\eta}{2}$, 则$|a-b| > \epsilon$, 矛盾. 故$a = b$.


\item[(4)]
\begin{itemize}
\item[(4.1)] $S = \{x| x^2 <2 \} = (-\sqrt 2, \sqrt 2)$, 上、下确界分别为$\sqrt 2$和$-\sqrt 2$. 只证上确界为$\sqrt 2$. 显然$\sqrt 2$为$S$的上界. 对任意的$\epsilon >0$, 不妨设$\epsilon < 2 \sqrt 2$. 取$x_0 = \sqrt 2 - \frac{\epsilon}{2}$, 则
\[
x_0^2 = (\sqrt 2 - \frac{\epsilon}{2})^2 = 2 + \frac{\epsilon^2}{4} - \sqrt 2 \epsilon < 2,
\]
即$x_0 \in S$, 且$x_0 > \sqrt 2 -\epsilon$. 因此$\sqrt 2$是$S$的上确界.

\item[(4.2)] $S = \{ x | x = n!, n \in \mathbb N^*\} = \{1, 2, 6, \cdots \}$上、下确界分别为$+\infty$和1. 首先1是$S$的下界, 且任何大于1的数都不是$S$的下界, 因此1是$S$的最大下界, 即下确界. 对任意的$M>0$, 取$n = [M] +1 \in \mathbb N^*$, 则$x = n! \geq n >M$, 故$S$无上界, 即$S$的上确界为$+\infty$.

\item[(4.3)] $S$的上、下确界分别为1和0. 只证上确界为1. 对任意的$\epsilon>0$, 不妨设$\epsilon<1$, 则$0<1-\epsilon<1$. 根据无理数的稠密性, 存在无理数$x_0 \in (1-\epsilon, 1)$. 于是$x_0 \in S$且$x_0 >1-\epsilon$. 即证.

\item[(4.4)] $S$的上、下确界分别为1和$\frac{1}{2}$. 因为$S$中的最小元素为$\frac{1}{2}$, 故$\frac{1}{2}$是$S$的最大下界, 即下确界. 由于$1 - \frac{1}{2^n}<1 (n \in \mathbb N^*)$, 所以1是$S$的上界. 对任意的$\epsilon >0$, 存在$n_0 \in \mathbb N^*$使得$\frac{1}{2^{n_0}} < \epsilon$, 于是取$x_0 = 1 -\frac{1}{2^{n_0}} \in S $, 且满足$x_0 > 1 - \epsilon$, 因此上确界为1.
\end{itemize}

\item[(5)] 考虑集合$\{ n \in \mathbb Z|\frac{x}{h} < n \}$是整数集的非空下有界子集, 根据确界原理存在下确界$k$, 即$(k-1) \leq \frac{x}{h} < k$. 由于下确界的唯一性, $k$也是唯一的.

\begin{itemize}
\item[(5.1)] 取$x=1, h = \epsilon$, 根据阿基米德原理存在$n \in \mathbb Z$使得$1 < \epsilon \cdot n$. 由于$\epsilon >0$, 故$n >0$, 即$n \in \mathbb N^*$.

\item[(5.2)] 利用上一结论反证.

\item[(5.3)] 利用(5.1), 存在$n \in \mathbb N^*$使得$0 < \frac{1}{n} < b-a$. 再根据阿基米德原理存在唯一的$m\in \mathbb Z$使得$\frac{m-1}{n} \leq a < \frac{m}{n}$. 于是$\frac{m}{n}<b$.

\item[(5.4)] 直接取$h=1$即得. 


\end{itemize}

\end{itemize}

\clearpage

\section{10月7--8日这一周}

\subsection*{题目}

\begin{itemize}

\item[1.] 证明以下不等式成立:
\[
|\sin x| \leq |x|, \ \forall x \in \mathbb R.
\]

\item[2.] 利用二项式定理证明算术-几何均值不等式: 设$a_1, a_2, \ldots, a_n$为非负实数, 则
\[ \frac{a_1 + a_2 + \cdots + a_n}{n} \geq \sqrt[n]{a_1a_2 \cdots a_n}. \eqno(1)
\]
其中等号成立的充分必要条件是$a_1 = a_2 = \cdots = a_n$.

\item[3.] 按照提示证明Cauchy-Schwarz不等式:
\[ \left( \sum_{k=1}^{n}a_k b_k\right)^2 \leq \left( \sum_{k=1}^{n}a_k^2\right) \left( \sum_{k=1}^{n}b_k^2\right).
\]
其中等号成立的条件: 存在$\lambda\in\mathbb R$使得$a_i=\lambda b_i$（$1\leq i\leq n$），或者$b_1=\cdots=b_n=0$.

\begin{itemize}
\item[3.1] 利用Young不等式或$|AB| \leq \frac{A^2+B^2}{2}$. tip: 作以下变形：
\[
\frac{ \sum\limits_{k=1}^n a_k b_k }{ \left( \sum\limits_{k=1}^n a_k^2 \right)^{\frac{1}{2}}  \left( \sum\limits_{k=1}^n b_k^2 \right)^{\frac{1}{2}} } =\sum\limits_{k=1}^n  \frac{a_k b_k }{ \left( \sum\limits_{k=1}^n a_k^2 \right)^{\frac{1}{2}}  \left( \sum\limits_{k=1}^n b_k^2 \right)^{\frac{1}{2}} } \leq 1.
\]

\item[3.2] 利用Lagrange恒等式
\[
\left(\sum\limits_{k=1}^n a_k^2 \right) \left( \sum\limits_{k=1}^n b_k^2\right) - \left( \sum\limits_{k=1}^n |a_k b_k|  \right)^2 = \frac{1}{2} \sum\limits_{k=1}^n \sum\limits_{i=1}^n (|a_k||b_i| - |a_i||b_k|)^2 =  \sum\limits_{k=1}^{n-1} \sum\limits_{j=k+1}^n (|a_k||b_j| - |a_j||b_k|)^2;
\]

\end{itemize}

\item[4] 证明当$a, b, c$为非负实数时, 成立不等式
\[
\sqrt[3]{abc} \leq \sqrt{\frac{ab + bc + ca}{3}} \leq \frac{a+b+c}{3}.
\]
tip: 利用均值不等式.

\end{itemize}

\clearpage

\subsection*{原稿答案}

\noindent\emph{以下为按题干归入本周的原稿部分参考答案。}

\begin{itemize}

\item 利用二项式定理证明均值不等式: 首先将$k+1$个正数$a_1, a_2, \ldots, a_{k+1}$的算术平均值分解如下:
\[
\frac{a_1 + a_2 + \cdots + a_{k+1}}{k+1} = \frac{a_1 + a_2 + \cdots + a_{k}}{k} + \frac{ka_{k+1} - (a_1 + a_2 + \cdots + a_{k})}{k(k+1)} = A + B. \eqno{(1)}
\]
将$a_1, a_2, \ldots, a_{k+1}$根据需要重新编号, 使得$a_{k+1}$是其中的最大数. 进而再分解式(1). 这个分解中的B项一定是非负数, 从而满足条件$A>0, B \geq 0$. 从二项式展开定理就有$(A+B)^{k+1} \geq A^{k+1} + (k+1)A^kB$, 其后证明不变.

\item 利用Young不等式证明C-S不等式: 在以下过程中取$p=q=2$即可.
\begin{align*}
\frac{ \sum\limits_{k=1}^n a_k b_k }{ \left( \sum\limits_{k=1}^n a_k^p \right)^{\frac{1}{p}}  \left( \sum\limits_{k=1}^n b_k^q \right)^{\frac{1}{q}} } & =  \sum\limits_{k=1}^n \frac{  a_k b_k }{ \left( \sum\limits_{k=1}^n a_k^p \right)^{\frac{1}{p}}  \left( \sum\limits_{k=1}^n b_k^q \right)^{\frac{1}{q}} } \\
& = \sum\limits_{k=1}^n \left( \frac{  a_k^p }{ \sum\limits_{k=1}^n a_k^p } \right)^{\frac{1}{p}} \left( \frac{  b_k^q }{ \sum\limits_{k=1}^n b_k^q } \right)^{\frac{1}{q}} (= a \cdot b) \\
& \leq \sum\limits_{k=1}^n \left[   \frac{1}{p}    \left( \frac{  a_k^p }{ \sum\limits_{k=1}^n a_k^p } \right)   +  \frac{1}{q}  \left( \frac{  b_k^q }{ \sum\limits_{k=1}^n b_k^q } \right)     \right] (= \frac{a^p}{p} + \frac{b^q}{q})\\
& = \frac{1}{p} \sum\limits_{k=1}^n \left( \frac{  a_k^p }{ \sum\limits_{k=1}^n a_k^p } \right) + \frac{1}{q} \sum\limits_{k=1}^n \left( \frac{  b_k^q }{ \sum\limits_{k=1}^n b_k^q } \right) =1.
\end{align*}
等号成立的条件用归纳法.

\item Lagrange恒等式的证明: 将等式左边表示成行列式的形式.
\[
\left | \begin{matrix}
\sum_{i=1}^n a_i^2 & \sum_{i=1}^n a_i b_i  \\
\sum_{i=1}^n a_i b_i & \sum_{i=1}^n b_i^2 \\
\end{matrix} \right |
\]
行列式对应的矩阵为
\[
\left ( \begin{matrix}
a_1 & a_2 & \cdots & a_n \\
b_1 & b_2 & \cdots & b_n \\
\end{matrix} \right ) 
\left ( \begin{matrix}
a_1 & b_1 \\
a_2 & b_2 \\
\vdots & \vdots \\
a_n & b_n
\end{matrix} \right )
\]
用Cauchy-Binet公式得
\[
\left | \begin{matrix}
\sum_{i=1}^n a_i^2 & \sum_{i=1}^n a_i b_i  \\
\sum_{i=1}^n a_i b_i & \sum_{i=1}^n b_i^2 \\
\end{matrix} \right |
= \sum_{1 \leq i < j \leq n} \left ( \begin{matrix}
a_i & a_j  \\
b_i & b_j  \\
\end{matrix} \right ) \left ( \begin{matrix}
a_i & b_j  \\
a_j & b_j  \\
\end{matrix} \right )  =  \sum_{1 \leq i < j \leq n}(a_i b_j - a_j b_i)^2.
\]

\item 证明当$a, b, c$为非负实数时, 成立不等式
\[
\sqrt[3]{abc} \leq \sqrt{\frac{ab + bc + ca}{3}} \leq \frac{a+b+c}{3}.
\]

\begin{proof}
左边的不等式利用均值不等式:
\[
\frac{ab+bc+ac}{3} \geq \sqrt[3]{a^2 b^2 c^2},
\]
再开平方即可. 右边的不等式两边平方后等价于
\[
(a+b+c)^2 \geq 3(ab+bc+ac).
\]
将左边平方展开后移项凑平方即可.
\end{proof}

\item 设$f : \mathbb R \to \mathbb R$满足方程
\[
f(x + y) = f(x) + f(y), x, y \in \mathbb R.
\]
证明: 对一切有理数$x$, 有$f(x) = xf(1)$.

\begin{proof}
\[
f(\frac{p}{q}) = f (\overbrace{\frac{1}{q} + \cdots + \frac{1}{q}}^{p\mbox{个}}) = p f(\frac{1}{q}) = p \frac{\overbrace{f(\frac{1}{q}) + \cdots +f(\frac{1}{q})}^{q\mbox{个}}}{q} = \frac{p}{q} f(1).
\]
\end{proof}

\item 给定函数$f : \mathbb R \to \mathbb R$, 如果$x \in \mathbb R$使得$f(x) =x$, 称$x$为$f(x)$的一个不动点. 若$f \circ f$有唯一的不动点, 证明$f$也有唯一的不动点.

\begin{proof}
先证明$f$的不动点一定是$f^2$的不动点. 再反证, 假设$f$没有不动点, 那么一定有$f(x) < x$或$f(x) >x$成立. 假设前者, 那么$f^2 (x) < f(x) < x, \forall x$成立, 说明$f^2$没有不动点, 矛盾.
\end{proof}


\end{itemize}

\clearpage

\section{10月12--13日这一周}

\subsection*{题目}

\begin{question}
设$A, B$是两个非空数集, 且$\forall a \in A, \forall b \in B$, 都有$a \leq b$. 证明: $\sup A \leq \inf B$.
\end{question}

\begin{question}
设$f : \mathbb R \to \mathbb R$满足方程
\[
f(x + y) = f(x) + f(y), x, y \in \mathbb R.
\]
证明: 对一切有理数$x$, 有$f(x) = xf(1)$.
\end{question}

Tips: 整数形式$\rightarrow$$1/n$形式$\rightarrow$$p/q$形式.

\begin{question}
函数
\[
\cosh x = \frac{e^x + e^{-x}}{2}, \ \sinh x = \frac{e^x - e^{-x}}{2},
\]
分别称为双曲余弦和双曲正弦. 分别求反函数.
\end{question}

\begin{question}
给定函数$f : \mathbb R \to \mathbb R$, 如果$x \in \mathbb R$使得$f(x) =x$, 称$x$为$f(x)$的一个不动点. 若$f \circ f$有唯一的不动点, 证明$f$也有唯一的不动点.
\end{question}

\begin{question}
平面直角坐标系中, 两坐标$x, y$均为有理数的点$(x, y)$称为有理点. 证明: 平面上全体有理点所成的集合为可数集.
\end{question}

\begin{question}
试证明代数无理数的全体是可数集. (选做)

\end{question}

\clearpage

\subsection*{原稿答案}

\noindent\emph{原稿未附这一周对应的参考答案。}

\clearpage

\section{10月19--20日这一周}

\subsection*{题目}

\begin{itemize}
\item[1.] 设$a_n>0 (n = 1, 2, 3, \ldots)$, 且$\lim_{n \to \infty} a_n =a$. 证明:
\[
\lim_{n \to \infty} \sqrt[n]{a_1 a_2 \cdots a_n} = a.
\]
\item[2.] 利用第一问的结论, 证明:
\begin{itemize}
\item[(1)] 若$a_n>0 (n = 1, 2, 3, \ldots)$, 且$\lim_{n \to \infty}\frac{a_{n+1}}{a_n} =a$, 则$\lim_{n \to \infty} \sqrt[n]{a_n} =a$;
\item[(2)] 当$a>0$时, $\lim_{n \to \infty} \sqrt[n]{a} = 1$; (请给出新的证明)
\item[(3)] $\lim_{n \to \infty} \sqrt[n]{n} =1$; (请给出新的证明)
\item[(4)] $\lim_{n \to \infty} (n!)^{-\frac{1}{n}} =0$.
\end{itemize} 
\end{itemize}

\begin{itemize}
\item[3.]
\begin{itemize}
\item[(1)] 设$a_n \leq b \leq c_n, n \in \mathbb N^*$, 且$\lim_{n \to \infty} (a_n - c_n) =0$. 证明:
\[
\lim_{n \to \infty} a_n = \lim_{n \to \infty} c_n = b.
\]
\item[(2)] 设$a_n \leq b_n \leq c_n, n \in \mathbb N^*$, 若已知$\lim_{n \to \infty}(c_n -a_n) =0$, 问: 数列$\{b_n\}$是否收敛?
\item[(3)] 设已知$\{a_n\}$收敛到$a$, 又有$b \leq a \leq c$, 问: 是否存在$N$, 使得当$n >N$时成立$b \leq a_n \leq c$?
\end{itemize}
\end{itemize}

\begin{itemize}
\item[4.]
如果$a_0 + a_1 + \cdots + a_p =0$, 求证
\[
\lim_{n \to \infty} (a_0 \sqrt n + a_1 \sqrt{n+1} + \cdots + a_p \sqrt{n+p})=0.
\]
tips: 利用已知结果: $\lim_{n \to \infty} (\sqrt{n+1} - \sqrt n) =0$.

\item[5.]
设数列$\{a_n\}$满足
\[
\lim_{n \to \infty} a_{2n-1} =a, \lim_{n \to \infty} a_{2n} =b.
\]
证明:
\[
\lim_{n \to \infty} \frac{a_1 + a_2 + \cdots + a_n}{n} = \frac{a+b}{2}.
\]

\item[6.]
证明:
\[
\lim_{n \to \infty} \sum_{k=1}^n \left( \sqrt{1 + \frac{k}{n^2} } -1\right) = \frac{1}{4}.
\]
tips: 利用平方差公式和迫敛性. 

\item[7.] 设$x_1 = \frac{c}{2}, x_{n+1} = \frac{c}{2} + \frac{x_n^2}{2}, n \in \mathbb N^*$, 证明: 若$c>1$, 则数列$\{x_n \}$发散.

\end{itemize}

\clearpage

\subsection*{原稿答案}

\noindent\emph{以下为按题干归入本周的原稿部分参考答案。}

\begin{itemize}
\item[1.] 设$a_n>0 (n = 1, 2, 3, \ldots)$, 且$\lim_{n \to \infty} a_n =a$. 证明:
\[
\lim_{n \to \infty} \sqrt[n]{a_1 a_2 \cdots a_n} = a.
\]

\begin{proof}
分情况讨论. 当$a=0$时, 使用迫敛性以及不等式
\[
0 \leq \sqrt[n]{a_1 \cdots a_n} \leq \frac{a_1 + \cdots + a_n}{n}.
\]
当$a\neq 0$时, 使用迫敛性, 极限的运算法则以及不等式
\[
\frac{1}{\frac{\frac{1}{a_1} + \cdots + \frac{1}{a_n}}{n}}\leq \sqrt[n]{a_1 \cdots a_n} \leq \frac{a_1 + \cdots + a_n}{n}.
\]
\end{proof}

\item[2.] 利用第一问的结论, 证明:
\begin{itemize}
\item[(1)] 若$a_n>0 (n = 1, 2, 3, \ldots)$, 且$\lim_{n \to \infty}\frac{a_{n+1}}{a_n} =a$, 则$\lim_{n \to \infty} \sqrt[n]{a_n} =a$;
\item[(2)] 当$a>0$时, $\lim_{n \to \infty} \sqrt[n]{a} = 1$; (请给出新的证明)
\item[(3)] $\lim_{n \to \infty} \sqrt[n]{n} =1$; (请给出新的证明)
\item[(4)] $\lim_{n \to \infty} (n!)^{-\frac{1}{n}} =0$.
\end{itemize} 

\begin{proof} 做以下变形:

(1) $\sqrt[n]{a_n} = \sqrt[n]{a_1 \frac{a_2}{a_1} \frac{a_3}{a_2}\cdots \frac{a_n}{a_{n-1}}}$.

(2) $\sqrt[n]{a} = \sqrt[n]{a \cdot 1 \cdot 1 \cdots 1}$.

(3) $\sqrt[n]{n} = \sqrt[n]{\frac{1}{1} \cdot \frac{2}{1} \cdot \frac{3}{2} \cdots \frac{n}{n-1}}$.

(4) $\sqrt[n]{\frac{1}{n!}} = \sqrt[n]{\frac{1}{1} \frac{1}{2} \cdots \frac{1}{n}}$.


\end{proof}


\item[4.]
如果$a_0 + a_1 + \cdots + a_p =0$, 求证
\[
\lim_{n \to \infty} (a_0 \sqrt n + a_1 \sqrt{n+1} + \cdots + a_p \sqrt{n+p})=0.
\]

\begin{proof}

\[
a_0 \sqrt n + a_1 \sqrt{n+1} + \cdots + a_p \sqrt{n+p} = a_0 \sqrt n + a_1 \sqrt{n+1} + \cdots + a_p \sqrt{n+p}  - (a_0 + a_1 + \cdots + a_p) \sqrt{n} 
\]
\[
= a_1 (\sqrt{n+1}-\sqrt{n}) + \cdots + a_p (\sqrt{n+p} - \sqrt{n}) \leq \frac{a_1}{\sqrt{n+1} + \sqrt{n}} + \cdots + \frac{p a_p}{\sqrt{n+p} + \sqrt{n}} \to 0 (n \to \infty).
\]


\end{proof}


\item[6.]
证明:
\[
\lim_{n \to \infty} \sum_{k=1}^n \left( \sqrt{1 + \frac{k}{n^2} } -1\right) = \frac{1}{4}.
\]

\begin{proof}

分子有理化:
\[
\sqrt{1 + \frac{k}{n^2} } -1 = \frac{k}{n^2} \cdot \frac{1}{\sqrt{1 + \frac{k}{n^2} } +1} \leq \frac{k}{2n^2},
\]
此外, 
\[
\sqrt{1 + \frac{k}{n^2} } -1 = \frac{k}{n^2} \cdot \frac{1}{\sqrt{1 + \frac{k}{n^2} } +1} \geq \frac{k}{n^2} \cdot \frac{1}{\sqrt{1 + \frac{n}{n^2} } +1} = \frac{k}{n^2} \cdot \frac{1}{\sqrt{1 + \frac{1}{n} } +1} \]
\[
\geq \frac{k}{n} \cdot \frac{1}{n + \sqrt{n(n+1)}} \geq \frac{k}{n(2n+1)}.
\]


也可以使用均值不等式:
\[
1 + \frac{x}{2+x}= \frac{2}{\frac{1}{1+x}+1} \leq \sqrt{(1+x) \cdot 1} \leq \frac{1+x +1}{2} = 1 + \frac{x}{2} (x >-1)
\]
因此
\[
\frac{k}{2n^2 +n}\leq \frac{k}{2n^2 +k}= \frac{\frac{k}{n^2}}{2 + \frac{k}{n^2}} \leq \sqrt{1 + \frac{k}{n^2} } -1 \leq \frac{k}{2n^2}.
\]


\end{proof}

\item[7.] 设$x_
1 = \frac{c}{2}, x_{n+1} = \frac{c}{2} + \frac{x_n^2}{2}, n \in \mathbb N^*$, 证明: 若$c>1$, 则数列$\{x_n \}$发散.

\begin{proof}

反证. 假设$\lim_{n \to \infty}x_n = \xi$. 在$x_{n+1} = \frac{c}{2} + \frac{x_n^2}{2}$两边令$n \to \infty$, 利用极限的运算法则, 得到$\xi = \frac{c}{2} + \frac{\xi^2}{2}$. 注意方程$x^2-2x+c$的判别式$\Delta = 4-4c<0$恒成立, 因此没有实根, 说明极限不存在. 


也可以配方说明$x_{n+1} - x_n = \frac{x_n^2}{2} - x_n + \frac{c}{2}  = \frac{1}{2} (x_n -1)^2 + \frac{c-1}{2} \geq \frac{c-1}{2} >0$. 归纳说明$x_{n+1} \geq \frac{c-1}{2}n + \frac{c}{2} $趋于无穷.

\end{proof}


\hrule
\
\end{itemize}

\clearpage

\section{10月26--27日这一周}

\subsection*{题目}

3.1和3.2节课后布置的作业以及以下问题.

\begin{itemize}

\item[1.] 若在$[a, b]$中有两个数列$\{ x_n \}$与$\{ y_n \}$满足$x_n - y_n \to 0$, 则两个数列中能找到具有相同下标$n_k$的子列, 使得
\[
x_{n_k} \to \alpha, \ y_{n_k} \to \alpha.
\]

\item[2.] 用紧致性定理证明单调有界原理.

\item[3.] 求极限$\lim_{n \to \infty} \frac{2^n n!}{n^n}$.

\item[4.] 设$\{a_n\}$是有界数列, 并满足
\[
a_{n+1} \geq a_n - \frac{1}{2^n}, \ n \geq 1.
\]
证明: $\{a_n\}$收敛.

\item[5.] 用基本列证明数列$\{ \sin n \}$发散.

\item[6.]
\begin{itemize}
\item[(1)] 利用均值不等式证明数列$a_n = \left( 1 + \frac{1}{n}\right)^n$严格递增, 而$b_n = \left( 1 + \frac{1}{n}\right)^{n+1}$严格递减.
\item[(2)] 证明不等式
\[
\left( 1 + \frac{1}{n}\right)^n < e < \left( 1 + \frac{1}{n}\right)^{n+1}, n \in \mathbb N^*.
\]
\item[(3)] 利用对数函数$\ln x$的严格递增性质证明
\[
\frac{1}{n+1} < \ln(1 + \frac{1}{n}) < \frac{1}{n}, n \in \mathbb N^*.
\]
\item[(4)] 对$n \in \mathbb N^*$, 求证
\[
\frac{1}{2} + \frac{1}{3} + \cdots + \frac{1}{n+1} < \ln(n+1) < 1 + \frac{1}{2} + \cdots + \frac{1}{n}.
\]
\item[(5)] 令
\[
x_n = 1 + \frac{1}{2} + \cdots +\frac{1}{n} - \ln n, n \in \mathbb N^*.
\]
证明: $\lim_{n \to \infty} x_n$存在. 这个极限常记为$\gamma$, 称为Euler常数.
\item[(6)] 求极限
\[
\lim_{n \to \infty} \left( \frac{1}{n+1} + \frac{1}{n+2} + \cdots + \frac{1}{n+n}\right).
\]
\item[(7)] 证明不等式
\[
\left( \frac{n+1}{e} \right)^n < n! < e \left( \frac{n+1}{e} \right)^{n+1}.
\]
\end{itemize}

\end{itemize}

\clearpage

\subsection*{原稿答案}

\noindent\emph{以下为按题干归入本周的原稿部分参考答案。}

\begin{itemize}
\item[2.] 用紧致性定理证明单调有界原理.


\begin{proof}
紧致性定理(或称列紧性定理, Bolzano-Weierstrass定理)即课本定理2.10致密性定理: 有界数列存在收敛子列. 


设$\{a_n \}$为单调有界数列, 下证数列$\{a_n\}$收敛. 首先根据紧致性定理, 存在子列$a_{n_k}$收敛$\xi$ ($k \to \infty$). 如此这个命题转化为证明: {\bf 如果单调数列有子列收敛, 那么该数列一定收敛.} 不妨设$\{a_n \}$单调增. 首先注意$\xi$是子列$a_{n_k}$的上界, 也一定是数列$\{a_n\}$的上界. 事实上若不然, 存在某个$a_m > \xi$. 那么根据单调性,  在子列$a_{n_k}$中的$a_{n_m} \geq a_m > \xi$. 从而有$\lim_{k \to \infty}a_{n_k} \geq a_{n_m} > \xi$, 矛盾.

其次根据极限的定义, $\forall \epsilon>0, \exists k_0 >0$, s.t. 当$k > k_0$时$|a_{n_k} - \xi| < \epsilon$. 由于$\{a_n \}$单调增, 那么当$n \geq n_{k_0}$时, $|a_n - \xi| \leq |a_{n_{k_0}} - \xi|<\epsilon$. 即证.

\end{proof}


\item[3.] 求极限$\lim_{n \to \infty} \frac{2^n n!}{n^n}$.


\begin{proof}
令$a_n = \frac{2^n n!}{n^n}$, 则$\frac{a_{n+1}}{a_n} = \frac{2}{(1+\frac{1}{n})^n} <1$当$n \geq 2$时成立. 因此$\{a_n\}$单调递减有下界0. 在等式
\[
a_{n+1} = a_n \frac{2}{(1+\frac{1}{n})^n}
\]
两边令$n \to \infty$, 得到极限为0.
\end{proof}

\item[4.] 设$\{a_n\}$是有界数列, 并满足
\[
a_{n+1} \geq a_n - \frac{1}{2^n}, \ n \geq 1.
\]
证明: $\{a_n\}$收敛.

\begin{proof}

构造辅助数列$b_{n+1} = a_{n+1} - \frac{1}{2^n}$, 则$b_{n+1} \geq b_n$. 由于$\{a_n\}$有界, $\{b_n\}$一定有界. 因此$\{ b_n\}$收敛, 再得到数列$\{a_n\}$收敛.


\end{proof}


\item[6.]
\begin{itemize}
\item[(1)] 利用均值不等式证明数列$a_n = \left( 1 + \frac{1}{n}\right)^n$严格递增, 而$b_n = \left( 1 + \frac{1}{n}\right)^{n+1}$严格递减.
\item[(2)] 证明不等式
\[
\left( 1 + \frac{1}{n}\right)^n < e < \left( 1 + \frac{1}{n}\right)^{n+1}, n \in \mathbb N^*.
\]
\item[(3)] 利用对数函数$\ln x$的严格递增性质证明
\[
\frac{1}{n+1} < \ln(1 + \frac{1}{n}) < \frac{1}{n}, n \in \mathbb N^*.
\]
\item[(4)] 对$n \in \mathbb N^*$, 求证
\[
\frac{1}{2} + \frac{1}{3} + \cdots + \frac{1}{n+1} < \ln(n+1) < 1 + \frac{1}{2} + \cdots + \frac{1}{n}.
\]
\item[(5)] 令
\[
x_n = 1 + \frac{1}{2} + \cdots +\frac{1}{n} - \ln n, n \in \mathbb N^*.
\]
证明: $\lim_{n \to \infty} x_n$存在. 这个极限常记为$\gamma$, 称为Euler常数.
\item[(6)] 求极限
\[
\lim_{n \to \infty} \left( \frac{1}{n+1} + \frac{1}{n+2} + \cdots + \frac{1}{n+n}\right).
\]
\item[(7)] 证明不等式
\[
\left( \frac{n+1}{e} \right)^n < n! < e \left( \frac{n+1}{e} \right)^{n+1}.
\]
\end{itemize}


\begin{proof}

(1) 注意数列$\{ a_n \}$的通项为一个数自乘$n$次, 为了利用均值不等式, 再乘上1可得
\[
a_n = 1 \cdot  \left( 1 + \frac{1}{n} \right)^n < \Big[ \frac{n \left( 1 + \frac{1}{n} \right) + 1 }{n+1} \Big]^{n+1} =  \left( \frac{ n+ 2}{n+1} \right)^{n+1} = a_{n+1}.
\]
利用同样的方法, 可以证明数列$b_n =  \left( 1 + \frac{1}{n} \right)^{n+1}$严格单调递减:
\[
\frac{1}{b_n} = 1 \cdot  \left( \frac{n}{n+1} \right)^{n+1} < \Big[ \frac{(n+1)\left( \frac{n}{n+1}\right)+1}{n+2}\Big]^{n+2} =  \left( \frac{n+1}{n+2} \right)^{n+2}  = \frac{1}{b_{n+1}}.
\] 
又由于对每个$n$成立不等式$a_n < b_n$, 可见$\{a_n \}$严格单调增加, 且以每一个$b_n$为上界; 同时$\{b_n \}$严格单调减少, 且以每一个$a_n$为下界. 因此两个数列都收敛, 利用$a_n (1 + \frac{1}{n}) =b_n$, 令$n \to \infty$可得极限相同.

(2) 利用(1)易得.

(3) 对(2)取对数易得.

(4) 注意利用对数函数的性质:
\[
\ln(n+1) = \ln 2 + \ln (1 + \frac{1}{2}) + \ln (1 + \frac{1}{3}) + \cdots + \ln (1 + \frac{1}{n}).
\]
在利用结论(3)易得.

(5)利用不等式
\[
\frac{1}{n+1} < \ln(1 + \frac{1}{n}) < \frac{1}{n}.
\]
考虑数列前后两项之差, 则
\[
x_{n+1} - x_n  = \frac{1}{n+1} - \ln(n+1) + \ln n = \frac{1}{n+1} - \ln \left( 1 + \frac{1}{n} \right) < 0.
\]
因此$\{x_n\}$严格单调递减. 以下只需证明数列有下界即可. 实际上
\[
1 + \frac{1}{2} + \cdots + \frac{1}{n} > \ln(n+1) = \ln n + \ln  \left( 1 + \frac{1}{n} \right) > \ln n + \frac{1}{n+1}.
\]
因此
\[
x_n > \frac{1}{n+1} >0.
\]
即证.


在确立了数列$\{ x_n \}$单调递减之后, 可以用和自然底数$e$时同样的方式引入第二个数列$\{d_n\}$, 其中$d_n = 1 + \frac{1}{2} + \cdots + \frac{1}{n} - \ln(n+1)$. 此时$d_n < x_n$, 且$x_n - d_n \to 0 (n \to \infty)$. 同样可以得到
\[
d_{n+1} - d_n = \frac{1}{n+1} - \ln(n+2) + \ln(n+1) = \frac{1}{n+1} - \ln \left( 1 + \frac{1}{n+1} \right) >0.
\]
因此$\{d_n\}$严格递增, 且$d_1 < d_2 < \cdots < d_n < x_n < \cdots < x_2 < x_1$. 因此$\{d_n\}$和$\{x_n\}$收敛到同一极限.

(6) 设$a_n = \frac{1}{n+1} + \frac{1}{n+2} + \cdots + \frac{1}{2n}$. 利用单调有界原理证明收敛. 比较相邻两项之差
\[
a_{n+1} - a_n =\frac{1}{2n+1} + \frac{1}{2n+2} - \frac{1}{n+1} >  2 \cdot \frac{1}{2n+2} - \frac{1}{n+1} =0.
\]
可见数列严格单增. 由于$a_n < \frac{n}{n+1} <1$, 因此有上界1, 则数列收敛.


又注意利用(5)
\[
x_{2n} - x_n =  \frac{1}{n+1} + \frac{1}{n+2} + \cdots + \frac{1}{2n} - \ln 2,
\]
因此
\[
\lim_{n \to \infty} \left(  \frac{1}{n+1} + \frac{1}{n+2} + \cdots + \frac{1}{2n} \right) = \lim_{n \to \infty}(x_{2n} - x_n) + \ln 2 = \ln 2.
\]


(7) 首先$\left( \frac{n+1}{e} \right)^n < n!  \Leftrightarrow \frac{(n+1)^n}{n!} < e^n$. 由于$\left(1+ \frac{1}{n} \right)^n <e$, 连乘可得
\[
\left(1+ \frac{1}{1} \right)^1 \cdot \left(1+ \frac{1}{2} \right)^2 \cdots \left(1+ \frac{1}{n} \right)^n < e^n,
\]
注意上式左边可写为
\[
2 \cdot (\frac{3}{2})^2 \cdot (\frac{4}{3})^3 \cdots (\frac{n+1}{n})^n =  \frac{(n+1)^n}{n!}.
\]
即证.


类似可证$e^n < \frac{(n+1)^{n+1}}{n!}$. 由$e < (\frac{n+1}{n})^{n+1}$连乘可得
\[
e^n < 2^2 \cdot (\frac{3}{2})^3 \cdot (\frac{4}{3})^4 \cdots (\frac{n+1}{n})^{n+1} = \frac{(n+1)^{n+1}}{n!}.
\]



\end{proof}



\end{itemize}

\clearpage

\section{11月2--3日这一周}

\subsection*{题目}

第四节全部习题以及以下问题:

\begin{itemize}

\item[1.] 设$\lim_{x \to x_0}f(x) > a $. 证明: 当$x$足够靠近$x_0$, 但$x \neq x_0$时, 有$f(x) >a$.

\item[2.] 设$f$在$(-\infty, a)$上递增, 且有一数列$\{x_n\}$满足$x_n < a, \forall n \geq 1$, $x _n \to a$且使得
\[
\lim_{n \to \infty}f(x_n) =A.
\]
证明: $f(a-)=A$.

\item[3.] 确定常数$a, b$使得下列等式成立:

\begin{itemize}
\item[(1)] $\lim_{x \to \infty} \left( \frac{x^2 +1}{x+1} - ax -b \right)=0$;
\item[(2)] $\lim_{x \to +\infty}(\sqrt{x^2-x+1} -ax -b)=0$;
\item[(3)] $\lim_{x \to -\infty}(\sqrt{x^2-x+1} -ax -b)=0$.
\end{itemize}

\item[4.] 证明: $\lim_{x \to \infty}(\sin \sqrt{x+1} - \sin \sqrt{x-1}) =0$.

\item[5.] 求极限: $\lim_{n \to \infty} \sin(\pi \sqrt{n^2+1})$.

\item[6.] 设$f(x), g(x)$是$\mathbb R$上的周期函数, 若$\lim_{x \to +\infty}[f(x) -g(x)] =0$, 则$f(x) = g(x)$恒成立.

\end{itemize}

\clearpage

\subsection*{原稿答案}

\noindent\emph{以下为按题干归入本周的原稿部分参考答案。}

\begin{itemize}
\item[1.] 设$\lim_{x \to x_0}f(x) > a $. 证明: 当$x$足够靠近$x_0$, 但$x \neq x_0$时, 有$f(x) >a$.

\begin{proof}
设$\lim_{x \to x_0}f(x) = A$, 利用极限的保号性, 取$\epsilon = \frac{A - a}{2}$, 那么当$x$充分接近$x_0$且$x \neq x_0$时成立$f(x) \geq A- \epsilon > a$.

\end{proof}

\item[2.] 设$f$在$(-\infty, a)$上递增, 且有一数列$\{x_n\}$满足$x_n < a, \forall n \geq 1$, $x _n \to a$且使得
\[
\lim_{n \to \infty}f(x_n) =A.
\]
证明: $f(a-)=A$.

\begin{proof}
根据单调有界原理, $f(a-)$为有限数或是$+\infty$. 只需证明$f(a-)$不能发散到$+\infty$. 不妨设$\{x_n\}$单调递增收敛到$a$. 取$\delta = a-x_0$, 那么$\forall y \in (a-\delta, a)$, 存在$k$使得$y \leq x_k$. 由于$f(x)$单调, 那么$f(y) \leq f(x_k)$. 令$k \to \infty$, 则$f(y) \leq A$. 这说明$f$在$(a-\delta, a)$上有界, 因此$f(a-)$也是有限数, 故只能是$A$.

\end{proof}

\item[3.] 确定常数$a, b$使得下列等式成立:

\begin{itemize}
\item[(1)] $\lim_{x \to \infty} \left( \frac{x^2 +1}{x+1} - ax -b \right)=0$;
\item[(2)] $\lim_{x \to +\infty}(\sqrt{x^2-x+1} -ax -b)=0$;
\item[(3)] $\lim_{x \to -\infty}(\sqrt{x^2-x+1} -ax -b)=0$.
\end{itemize}


\begin{proof}
只证(1). 证法有二.

证一. 通分可得
\[
\lim_{x \to \infty} \frac{(1-a)x^2 - (a+b)x +1-b}{x+1}=0.
\] 
那么必有$1-a=0, a+b=0$. 从而$a=1, b=-1$.


证二. 将已知等式变形
\[
\lim_{x \to \infty}x \left( \frac{x^2+1}{x(x+1)} -a\right) =b,
\]
则必有
\[
\lim_{x \to \infty} \left( \frac{x^2+1}{x(x+1)} -a\right) =0,
\]
即
\[
a = \lim_{x \to \infty}  \frac{x^2+1}{x(x+1)} =1.
\]
于是
\[
b = \lim_{x \to \infty} \left( \frac{x^2+1}{x+1} -ax\right) = \lim_{x \to \infty} \frac{1-x}{x+1} = -1.
\]

\end{proof}

\item[4.] 证明: $\lim_{x \to \infty}(\sin \sqrt{x+1} - \sin \sqrt{x-1}) =0$.

\begin{proof}
和差化积即可,
\[
\sin \sqrt{x+1} - \sin \sqrt{x-1} = 2 \sin \frac{\sqrt{x+1} - \sqrt{x-1}}{2} \cdot \cos  \frac{\sqrt{x+1} + \sqrt{x-1}}{2}
\]
\[
= 2 \sin \frac{1}{\sqrt{x+1} + \sqrt{x-1}} \cdot \cos  \frac{\sqrt{x+1} + \sqrt{x-1}}{2}.
\]

\end{proof}

\item[5.] 求极限: $\lim_{n \to \infty} \sin(\pi \sqrt{n^2+1})$.


\begin{proof}
做变形
\[
\sin((\sqrt{n^2+1} -n)\pi + n \pi )= (-1)^n \sin (\sqrt{n^2+1} -n)\pi = (-1)^n \sin \frac{\pi}{\sqrt{n^2+1} +n}.
\]

\end{proof}


\item[6.] 设$f(x), g(x)$是$\mathbb R$上的周期函数, 若$\lim_{x \to +\infty}[f(x) -g(x)] =0$, 则$f(x) = g(x)$恒成立.

\begin{proof}
设$f(x)$的周期为$T$. 根据题设$\lim_{x \to +\infty}[f(x+T) -g(x+T)] =0$, 从而由不等式
\[
|g(x+T) - g(x) | \leq |g(x+T) - f(x+T)| + |f(x+T) - f(x)| + |f(x) - g(x)|
\]
可得$g(x+T) - g(x) \to 0 ( x \to +\infty)$. 因为$g(x+T) - g(x)$为周期函数, 则$g(x+T) - g(x) \equiv 0$. 即$g(x+T) = g(x)$对所有$x \in \mathbb R$成立. 这说明$g(x)$也是周期为$T$的函数, 即$f(x) - g(x)$周期为$T$. 由条件$\lim_{x \to +\infty}[f(x) -g(x)] =0$即证.

\end{proof}

\end{itemize}

\clearpage

\section{11月9--10日这一周}

\subsection*{题目}

课后习题3.5节1,2,4,5以及4.1节1,2,3题.

\begin{itemize}

\item[(1)] 求极限: 
\[
1. \ \lim_{x \to 0} \frac{\sqrt{1 + x \sin x} -1}{x \arctan x}
\]

\[
2. \ \lim_{x \to 0} \frac{\sin 2x + 2 \arctan 3x + 3x^2}{\ln(1 + 3x + \sin^2 x) + xe^x} 
\]
\[
3. \lim_{x \to +\infty} \left( \frac{e^x + e^{-x}}{e^x - e^{-x}}\right)^{e^{2x}} 
\]
\[
4. \ \lim_{x \to 0+}(\sin x)^{\frac{1}{\ln x}}
\]

\item[(2)] 设当$x \to x_0$时, $\alpha(x) \sim \beta(x)$是等价无穷小量. 证明

1. \ $e^{\alpha(x)} - e^{\beta(x)}$是比$\alpha(x)$高阶的无穷小量;

2. \ $\ln \alpha(x) \sim \ln \beta(x) $.

\item[(3)] 有时用函数在一个点的振幅来刻画连续性, 定义如下: 对于点$a$的邻域$U(a, \delta)$, 定义$f$在这个邻域上的振幅为
\[
\omega_f(a, \delta) = \sup_{x \in U(a, \delta)}\{ f(x)\} - \inf_{x \in U(a, \delta)}\{ f(x)\},
\]
然后令
\[
\omega_f(a) = \lim_{\delta \to 0+} \omega_f(a, \delta),
\]
称为函数$f$在点$a$的振幅. {\bf 证明}: 函数$f$在点$a$处连续的充分必要条件是$\omega_f(a) =0$.

\item[(4)] 设函数$f \in C([a, b])$. 若有数列$\{x_n\} \subset [a, b]$, 使得$\lim_{n \to \infty}f(x_n)=A$, 证明: 存在$\xi \in [a, b]$使得$f(\xi) = A$.

\item[(5)] 集合$E$被称为开集, 指它的每个点皆为内点, 即$\forall x_0 \in E, \exists \delta >0$使得$(x_0 - \delta, x_0 + \delta ) \subset E$. 开集的补集被称为闭集. 因此对于一维欧氏空间, 开集即是可数个开区间的并集.

{\bf 证明命题}: 函数$f(x)$在$\mathbb R$上连续$\Leftrightarrow$任何开集的逆像仍为开集(即: 设$E$为$\mathbb R$上的开集, 则$f^{-1}(E) := \{ x | f(x) \in E \}$为实轴上的开集).

\item[(6)] 
设有三个函数$f_1, f_2, f_3 \in C([a, b])$, 对每个$x \in [a, b]$, 定义$f(x)$是$f_1(x), f_2(x), f_3(x)$中处于中间的一个值, 证明: $f \in C([a, b])$.

\item[(7)] 设$|x| <1$, 求极限
\[
\lim_{n \to \infty} \left( 1+ \frac{x + x^2 + \cdots +x^n}{n} \right)^n.
\]

\end{itemize}

\clearpage

\subsection*{原稿答案}

\noindent\emph{以下为按题干归入本周的原稿部分参考答案。}

\begin{itemize}
\item[(1)] 求极限: 
\[
1. \ \lim_{x \to 0} \frac{\sqrt{1 + x \sin x} -1}{x \arctan x}
\]
\[
2. \ \lim_{x \to 0} \frac{\sin 2x + 2 \arctan 3x + 3x^2}{\ln(1 + 3x + \sin^2 x) + xe^x} 
\]
\[
3. \lim_{x \to +\infty} \left( \frac{e^x + e^{-x}}{e^x - e^{-x}}\right)^{e^{2x}} 
\]
\[
4. \ \lim_{x \to 0+}(\sin x)^{\frac{1}{\ln x}}
\]


\begin{proof}


\[
1. \lim_{x \to 0} \frac{\sqrt{1 + x \sin x} -1}{x \arctan x} = \lim_{x \to 0}\frac{\frac{1}{2}x\sin x}{x^2} = \frac{1}{2}.
\]

\[
2.  \lim_{x \to 0} \frac{\sin 2x + 2 \arctan 3x + 3x^2}{\ln(1 + 3x + \sin^2 x) + xe^x} = \lim_{x \to 0} \frac{\frac{\sin 2x + 2 \arctan 3x + 3x^2}{x} }{\frac{\ln(1 + 3x + \sin^2 x) + xe^x}{x}} = \lim_{x \to 0} \frac{2\frac{\sin 2x}{2x} + 6 \frac{\arctan 3x}{3x}}{(3 + \frac{\sin^2 x}{x})+1} =2.
\]


3. 作变量代换$t = e^{-2x}$, 则$x \to +\infty$时$t \to 0+$. 故
\[
 \lim_{x \to +\infty} \left( \frac{e^x + e^{-x}}{e^x - e^{-x}}\right)^{e^{2x}} = \lim_{t \to 0+}\left( \frac{1+t}{1-t}\right)^{\frac{1}{t}} = \lim_{t \to 0+}\left( 1 + \frac{2t}{1-t}\right)^{\frac{1}{t}} = e^2.
\]

4. 做指数化变形
\[
\lim_{x \to 0+}(\sin x)^{\frac{1}{\ln x}} = \lim_{x \to 0+} e^{\frac{\ln \sin x}{\ln x}} =  \lim_{x \to 0+} e^{\frac{\ln \frac{\sin x}{x} +\ln x}{\ln x}} =e.
\]

\end{proof}

\item[(2)] 设当$x \to x_0$时, $\alpha(x) \sim \beta(x)$是等价无穷小量. 证明

1. \ $e^{\alpha(x)} - e^{\beta(x)}$是比$\alpha(x)$高阶的无穷小量;

2. \ $\ln \alpha(x) \sim \ln \beta(x) $.

\begin{proof}
易证.
\end{proof}

\item[(3)] 有时用函数在一个点的振幅来刻画连续性, 定义如下: 对于点$a$的邻域$U(a, \delta)$, 定义$f$在这个邻域上的振幅为
\[
\omega_f(a, \delta) = \sup_{x \in U(a, \delta)}\{ f(x)\} - \inf_{x \in U(a, \delta)}\{ f(x)\},
\]
然后令
\[
\omega_f(a) = \lim_{\delta \to 0+} \omega_f(a, \delta),
\]
称为函数$f$在点$a$的振幅. {\bf 证明}: 函数$f$在点$a$处连续的充分必要条件是$\omega_f(a) =0$.

\begin{proof}

$\Rightarrow$. 利用连续性的定义, $\forall \epsilon>0$充分小, $\exists \delta>0$, 使得当$x \in U(a, \delta)$时, $|f(x) - f(x_0)| < \epsilon$. 则$\forall y_1, y_2 \in U(a, \delta)$, 则$|f(y_1) - f(y_2)|  < 2 \epsilon$. 注意$\sup_{x \in U(a, \delta)}\{ f(x)\} - \inf_{x \in U(a, \delta)}\{ f(x)\} = \sup_{x, x' \in U(a, \delta)} |f(x) - f(x')|$. 因此$\omega_f(a, \delta) < 2\epsilon$. 即证. 

$\Leftarrow$. 类似可证.


\end{proof}

\item[(4)] 设函数$f \in C([a, b])$. 若有数列$\{x_n\} \subset [a, b]$, 使得$\lim_{n \to \infty}f(x_n)=A$, 证明: 存在$\xi \in [a, b]$使得$f(\xi) = A$.

\begin{proof}
从$\{x_n\}$中取出收敛子列$x_{n_k} \to \xi$, 则$\xi \in [a, b]$. 根据连续性, 
\[
A = \lim_{n \to \infty}f(x_n) = \lim_{k \to \infty}f(x_{n_k}) = f( \lim_{k \to \infty} x_{n_k}) = f(\xi).
\]

\end{proof}


\item[(5)] 集合$E$被称为开集, 指它的每个点皆为内点, 即$\forall x_0 \in E, \exists \delta >0$使得$(x_0 - \delta, x_0 + \delta ) \subset E$. 开集的补集被称为闭集. 因此对于一维欧氏空间, 开集即是可数个开区间的并集.

{\bf 证明命题}: 函数$f(x)$在$\mathbb R$上连续$\Leftrightarrow$任何开集的逆像仍为开集(即: 设$E$为$\mathbb R$上的开集, 则$f^{-1}(E) := \{ x | f(x) \in E \}$为实轴上的开集).

\begin{proof}

$(\Rightarrow)$ 要证$f^{-1}(E)$为开集, 即要证明: $\forall x_0 \in f^{-1}(E), \exists \delta >0$使
\[
(x_0 - \delta, x_0 + \delta ) \subset f^{-1}(E).
\]
由$x_0 \in f^{-1}(E)$, 知$y_0 = f(x_0) \in E$. 既然$E$为开集, 所以$\exists \epsilon>0$, 使得
\[
(f(x_0) - \epsilon, f(x_0) + \epsilon) = (y_0 - \epsilon, y_0 + \epsilon) \subset E.
\]
由于$f$连续, 对$y_0$的邻域$(y_0 - \epsilon, y_0 + \epsilon)$, $\exists \delta >0$, 使得
\[
f((x_0 - \delta, x_0 + \delta )) \subset (y_0 - \epsilon, y_0 + \epsilon) \subset E,
\]
从而$(x_0 - \delta, x_0 + \delta ) \subset f^{-1}((y_0 - \epsilon, y_0 + \epsilon)) \subset f^{-1}(E)$, 故$f^{-1}(E)$为开集.


$(\Leftarrow)$ 已知任何开集的逆像仍为开集, 故$\forall x_0 \in \mathbb R, \forall \epsilon>0$, (设$y_0 = f(x_0)$) $(y_0 - \epsilon, y_0 + \epsilon)$的逆像$f^{-1}((y_0 - \epsilon, y_0 + \epsilon))$仍为开集. 由此对于$x_0 \in f^{-1}((y_0 - \epsilon, y_0 + \epsilon)), \exists \delta >0$, 使得$(x_0 - \delta, x_0 + \delta ) \subset f^{-1}((y_0 - \epsilon, y_0 + \epsilon))$. 故
\[
f((x_0 - \delta, x_0 + \delta )) \subset (y_0 - \epsilon, y_0 + \epsilon) = (f(x_0) - \epsilon, f(x_0) + \epsilon), 
\]
所以$f$连续.
\end{proof}


\item[(6)] 
设有三个函数$f_1, f_2, f_3 \in C([a, b])$, 对每个$x \in [a, b]$, 定义$f(x)$是$f_1(x), f_2(x), f_3(x)$中处于中间的一个值, 证明: $f \in C([a, b])$.

\begin{proof}

注意$f = f_1 + f_2 + f_3 - \max\{ f_1, f_2, f_3\} - \min\{ f_1, f_2, f_3\}$.

\end{proof}


\item[(7)] 设$|x| <1$, 求极限
\[
\lim_{n \to \infty} \left( 1+ \frac{x + x^2 + \cdots +x^n}{n} \right)^n.
\]


\begin{proof}


只需计算
\[
\lim_{n \to \infty} n\left( \frac{x + x^2 + \cdots +x^n}{n} \right) = \frac{x}{1-x}.
\]


\end{proof}


\end{itemize}

\clearpage

\section{11月16--17日这一周}

\subsection*{题目}

\begin{itemize}
\item[(1)] 利用二分法和紧致性证明有界性定理: 设$f \in C([a, b])$, 那么$f$在$[a, b]$上有界.

\item[(2)] 设$f$在$[a, b]$上连续, 若对每一个$x \in [a, b]$存在$y \in [a, b]$使得$|f(y)| \leq \frac{1}{2}|f(x)|$. 证明: $f$在$[a, b]$中有零点.

\item[(3)] 证明: 偶数次多项式$f(x) = x^{2n} + a_1 x^{2n-1} + \cdots + a_{2n-1} x + a_{2n} = 0$可以没有实根, 但当$a_{2n}<0$时则至少有两个实根.

\item[(4)] 设$f \in C((a, b))$, $f(a+)$和$f(b-)$有限. 证明: $f$在$(a, b)$上有界.

\item[(5)] 证明:

(1.) 设$f \in C([a, b])$, 若$(a, b)$中无$f(x)$的极值点, 则$f(x)$在$[a, b]$上单调.

(2.) 设$f \in C(\mathbb R)$. 若对任意的开区间$(a, b)$, 其值域$f((a, b))$必是开区间, 则$f(x)$必单调.

\item[(6)] 设$f \in C([a, b])$, $a \leq x_1 < x_2 < \cdots <x_n \leq b$. 证明: 存在$\xi \in [x_1, x_n]$, 使成立
\[
f(\xi) = \frac{f(x_1) + f(x_2) + \cdots + f(x_n)}{n}.
\]

\item[(7)] 若$f \in C([a, +\infty))$, 且存在有限极限$\lim_{x \to +\infty}f(x)$. 证明: $f$在$[a, +\infty)$上有界.
\end{itemize}

\clearpage

\subsection*{原稿答案}

\noindent\emph{以下为按题干归入本周的原稿部分参考答案。}

\begin{itemize}
\item[(1)] 利用二分法和紧致性证明有界性定理: 设$f \in C([a, b])$, 那么$f$在$[a, b]$上有界.

\begin{proof}

反证. 假设$f$在$[a, b]$上无界. 考察两个闭区间$[a, \frac{a+b}{2}]$和$[\frac{a+b}{2}, b]$, 可以断定$f$至少在其中一个上无界. 记这个区间为$[a_1, b_1]$. 重复上述讨论, 得到一串闭区间
\[
[a, b] \supset [a_1, b_1] \supset \cdots \supset [a_k, b_k] \supset \cdots
\]
满足:
\begin{itemize}
\item[(1)] $0 < b_k - a_k = \frac{1}{2^k}(b-a) $,
\item[(2)] $f$在$[a_k, b_k]$上无界.
\end{itemize}
令$\xi = \lim_{k \to \infty}a_k = \lim_{k \to \infty} b_k$. 因为函数$f$在$\xi$处连续, 所以存在$\eta>0$使得$f$在$(\xi -\eta, \xi + \eta)$上有界. 取$k$充分大使得$[a_k, b_k] \subset (\xi -\eta, \xi + \eta)$, 得到矛盾. 

\end{proof}



\item[(2)] 设$f$在$[a, b]$上连续, 若对每一个$x \in [a, b]$存在$y \in [a, b]$使得$|f(y)| \leq \frac{1}{2}|f(x)|$. 证明: $f$在$[a, b]$中有零点.


\begin{proof}[证一]
任取一点$x_0 \in [a, b]$, 那么存在点$x_1 \in [a, b]$, 使$|f(x_1)| \leq \frac{1}{2} |f(x_0)|$. 归纳进行下去得到一个数列$\{ x_n \}$满足
\[
|f(x_n)| \leq \frac{1}{2} |f(x_{n-1})| \leq \cdots \leq \frac{1}{2^n} |f(x_0)|.
\]
因此$\lim_{n \to \infty}f(x_n) =0$.


对数列使用紧致性定理, 存在收敛子列$\{ x_{n_k}\}$. 设其极限为
\[
\lim_{k \to \infty}x_{n_k} = \eta,
\]
则$\eta \in [a, b]$. 由于$f$在$\eta$处连续, 那么
\[
\lim_{k \to \infty}f(x_{n_k}) = f(\eta).
\]
由于$\{ f(x_{n_k})\}$是$\{ f(x_n)\}$的子列, 则$f(\eta)=0$.
\end{proof}


注意在构造数列$\{x_n\}$时, 可能某个$x_n$恰为$f$的零点, 则归纳过程结束.


\begin{proof}

注意此时$|f| \in C([a, b])$, 对$|f|$使用零点定理, 在题设条件下可见最小值只能是0. 这时最小值点是$|f|$的零点, 也是$f$的零点.
\end{proof}



\item[(3)] 证明: 偶数次多项式$f(x) = x^{2n} + a_1 x^{2n-1} + \cdots + a_{2n-1} x + a_{2n} = 0$可以没有实根, 但当$a_{2n}<0$时则至少有两个实根.


\begin{proof}

只需注意$\lim_{x \to -\infty} f(x) = \lim_{x \to +\infty}f(x) = +\infty$即可.

\end{proof}


\item[(4)] 设$f \in C((a, b))$, $f(a+)$和$f(b-)$有限. 证明: $f$在$(a, b)$上有界.

\begin{proof}[证明一]

在闭区间$[a, b]$上构造辅助函数
\[
F(x) = \begin{cases} 
f(x), & x \in (a, b)\\
f(a+), & x =a\\
f(b-), & x =b
\end{cases}
\]
那么$F \in C([a, b])$. 对$F$应用连续函数有界性定理, 得到$F$在$[a, b]$上有界, 因此$f$在$(a, b)$上有界.
\end{proof}


\begin{proof}[证明二]

根据函数极限的局部有界性定理, 存在$\delta>0$, 使得函数$f$在$(a, a+\delta)$和$(b-\delta, b)$上有界, 这里可取$a+\delta < b - \delta$. 然后在闭区间$[a+\delta, b-\delta]$上应用有界性定理即可.

\end{proof}


第二种证明方法可以处理$a=-\infty$和$b = +\infty$的情况.



\item[(5)] 证明:

(1.) 设$f \in C([a, b])$, 若$(a, b)$中无$f(x)$的极值点, 则$f(x)$在$[a, b]$上单调.

(2.) 设$f \in C(\mathbb R)$. 若对任意的开区间$(a, b)$, 其值域$f((a, b))$必是开区间, 则$f(x)$必单调.


\begin{proof}

(1) 注意最值一定是极值, 因此函数$f$在端点$a, b$处取到最值. 不妨设$f(a)$为最小值, $f(b)$为最大值. 反证, 设存在$x_1 < x_2$使得$f(x_1) > f(x_2)$. 那么在区间$[a, x_2]$上函数取到最大值, 且最大值非$f(x_2)$, 矛盾.

(2) 类似可证. 


\end{proof}


\item[(6)] 设$f \in C([a, b])$, $a \leq x_1 < x_2 < \cdots <x_n \leq b$. 证明: 存在$\xi \in [x_1, x_n]$, 使成立
\[
f(\xi) = \frac{f(x_1) + f(x_2) + \cdots + f(x_n)}{n}.
\]


\begin{proof}
利用最值定理和介值定理易得.

\end{proof}



\item[(7)] 若$f \in C([a, +\infty))$, 且存在有限极限$\lim_{x \to +\infty}f(x)$. 证明: $f$在$[a, +\infty)$上有界.
\end{itemize}


\begin{proof}

设$\lim_{x \to +\infty} f(x) = A$. 存在$X>0$充分大, 使得当$x > X$时$|f(x) | \leq A+1$. 以下易证.

\end{proof}

\clearpage

\section{12月7--8日这一周}

\subsection*{题目}

\begin{enumerate}
\item 设函数$f(x)$在点$a$处连续, 且$f(a) \neq 0, [f(x)]^2$在点$a$可导. 证明: $f(x)$也在点$a$可导. 
\item 求极限:
\begin{enumerate}
\item $\lim\limits_{x \to a} \frac{\sin x - \sin a}{\sin(x -a)}$.
\item $\lim\limits_{x \to 0} \frac{f(x) e^x -f(0)}{f(x) \cos x - f(0)}$ ($f'(0) \neq 0$).
\item $\lim\limits_{x \to x_0} \Big[ \frac{f(x)}{f(x_0)}\Big]^{\frac{1}{\ln x- \ln x_0}}$ ($x_0>0$, $f'(x_0)$存在, $f(x_0)>0$).
\item $\lim\limits_{x \to 0} \frac{1}{x} \sum_{i=1}^k f\left( \frac{x}{i}\right)$ (已知$f'(0)$存在, 且$f(0) =0$, $k \in \mathbb N^*$ ).
\end{enumerate}
\item 求和式$\sum_{k=1}^n \frac{1}{2^k} \tan \frac{x}{2^k} $.
\item 设$f(x)$为多项式, $a$是实数. 若$f(a) =0$, 则称$a$为$f(x)$的一个根. 若有因式分解$f(x) = (x - a)^k g(x)$, 其中$g(x)$也是多项式(注意常数是零次多项式), $g(a) \neq 0$, $k \geq 1$为正整数, 则称$a$是$f(x)$的$k$重根. 若$k=1$, 则称$x = a$为$f(x)$的单根. 

证明: 若$x =a $是$f(x)$的$k$重根, 则$x =a $也是$f'(x)$的$k-1$重根.
\item 设$p(x)$是有$n$个实根的$n$次多项式, 记它的相异根为$x_1, x_2, \cdots, x_k$, 其中根$x_i$的重数为$n_i, i=1, 2, \cdots, k, n_1 + n_2 + \cdots + n_k =n$. 证明:
\[
p'(x) = p(x) \sum_{i=1}^k \frac{n_i}{x-x_i}.
\]
\item 设$f(x) = x(x-1)(x-2) \cdots (x-100)$, 求$f'(0)$.
\item 设$y = \arctan x$, 求$y^{(n)}(0)$.
\item 设$y = f(x)$为单调可导函数, 其反函数是$x = \varphi(y)$, 已知$f(1) = 2, f'(1) = - \frac{1}{\sqrt 3}, f''(1) = 1$, 求$\varphi''(2)$.
\item 设函数$y = f(x)$在点$x$三次可导, 且$f'(x) \neq 0$. 若$f(x)$存在反函数$x = f^{-1}(y)$, 求$(f^{-1})'''(y)$.
\end{enumerate}

\clearpage

\subsection*{原稿答案}

\noindent\emph{以下为按题干归入本周的原稿部分参考答案。}

\begin{enumerate}
\item 求和式$\sum_{k=1}^n \frac{1}{2^k} \tan \frac{x}{2^k} $.

\begin{proof}
注意$\left( \ln \Big| \cos \frac{x}{2^k}\Big|\right)' = - \frac{1}{2^k} \tan \frac{x}{2^k} $. 因此
\[
\sum_{k=1}^n \frac{1}{2^k} \tan \frac{x}{2^k} = - \left( \sum_{k=1}^n  \ln \Big| \cos \frac{x}{2^k}\Big|\right)' = - \Big[ \ln \Big| \cos \frac{x}{2} \cdot \cos \frac{x}{2^2} \cdots \cos \frac{x}{2^n} \Big| \Big]'
\]
\[
= -\Big[ \ln \Big| \frac{\sin x}{2^n \sin \frac{x}{2^n}} \Big| \Big]' = \frac{1}{2^n} \cot \frac{x}{2^n} - \cot x.
\]
\end{proof}

\item 设函数$y = f(x)$在点$x$三次可导, 且$f'(x) \neq 0$. 若$f(x)$存在反函数$x = f^{-1}(y)$, 求$(f^{-1})'''(y)$.


\begin{proof}

已知$(f^{-1})'(y ) = \frac{1}{f'(x)}$, 以及$(f^{-1})''(y ) = - \frac{f''(x)}{(f'(x))^3}$. 利用链式法则
\[
(f^{-1})'''(y ) = ((f^{-1})''(y ))' = \left(- \frac{f''(x)}{(f'(x))^3} \right)' \cdot (f^{-1})'(y ) 
\]
\[
= - \frac{ f'''(x) \cdot (f'(x))^3 - f''(x) \cdot 3 (f'(x))^2 \cdot f''(x)        }{(f'(x))^6} \cdot \frac{1}{f'(x)} = \frac{3(f''(x))^2 - f'(x) f'''(x)}{(f'(x))^5}.
\]


\end{proof}
\end{enumerate}

\clearpage

\section{12月14--15日这一周}

\subsection*{题目}

\begin{enumerate}

\item 设$f(x)$在$[a, b]$上连续, 在$(a, b)$内可导. 求证: 若$f(x)$不是线性函数(即$ax+b$形式), 则在$(a, b)$内存在$x_1$和$x_2$, 使得
\[
f'(x_1) < \frac{f(b) - f(a)}{b-a} < f'(x_2).
\]

\item 设$f(x)$在$[a, b]$上二次可导, $f(a) = f(b) =0$. 若存在$a < c < b$使得$f(c) >0$, 则存在$\xi \in (a, b)$, 使得$f''(\xi) <0$.

\item 证明: 设$f'(x)$在$x=a$处的右极限存在且有限, 则$f(x)$在$x=a$的右极限也存在且有限.

\item 设函数$f(x)$在闭区间$[a, b]$上连续, 在开区间$(a, b)$内可导, 且$f(a) = f(b)$. 求证: 若$f(x)$在$[a, b]$上不为常数, 则$\exists \xi, \eta \in (a, b)$, 使得
\[
f'(\xi) \cdot f'(\eta) <0.
\]

\item 证明等式、不等式或求极限:
\begin{enumerate}
\item 当$0 < \alpha < \beta < \frac{\pi}{2}$时, $\frac{\beta - \alpha}{\cos^2 \alpha} < \tan \beta - \tan \alpha < \frac{\beta - \alpha}{\cos^2 \beta}$.

\item $\arctan e^x + \arctan e^{-x} = \frac{\pi}{2}, x \in (-\infty, +\infty)$.

\item 当$x>1$时, $e^x > e x$.

\item $\lim_{x \to a} \frac{\sin x^x - \sin a^x}{a^{x^x} - a^{a^x}} (a > 1)$.

\item $\lim_{n \to \infty} n[\arctan[\ln(n+1)] - \arctan (\ln n)]$.

\end{enumerate}

\item 函数$f(x)$在$[a, b]$可导, 其中$a >0$, 证明: 存在$\xi \in (a, b)$, 使得
\[
2 \xi [f(b) - f(a)] = (b^2 - a^2) f'(\xi).
\]

\item 设$f(x)$在$[0, 1]$上连续, 在$(0, 1)$内可导, 且$f(0) =0, f(1) = 1$, 证明:
\begin{enumerate}
\item 存在$\xi \in (0, 1)$使得$f(\xi) = 1 - \xi$.
\item 存在两个不同点$\eta_1, \eta_2 \in (0, 1)$, 使得$f'(\eta_1) f'(\eta_2) = 1$.
\end{enumerate}

\end{enumerate}

\clearpage

\subsection*{原稿答案}

\noindent\emph{以下为按题干归入本周的原稿部分参考答案。}

\begin{enumerate}
\item 设$f(x)$在$[a, b]$上连续, 在$(a, b)$内可导. 求证: 若$f(x)$不是线性函数(即$ax+b$形式), 则在$(a, b)$内存在$x_1$和$x_2$, 使得
\[
f'(x_1) < \frac{f(b) - f(a)}{b-a} < f'(x_2).
\]

\begin{proof}

过点$(a, f(a))$和$(b, f(b))$的线性函数即
\[
L(x) = f(a) + \frac{f(b) - f(a)}{b-a} (x-a).
\]
由于$f(x) \neq L(x)$, 故$\exists c \in (a, b)$使得点$(c, f(c))$在$L(x)$上方或下方. 不妨设点$(c, f(c))$在$L(x)$下方, 即
\[
f(c) < L(c).
\]
于是
\[
 \frac{f(c) - f(a)}{c-a} <  \frac{f(b) - f(a)}{b-a},  \frac{f(c) - f(b)}{c-b} >  \frac{f(b) - f(a)}{b-a}.
\]
进一步在$[a, c]$和$[c, b]$上应用Lagrange中值定理, 即证.

\end{proof}



\item 证明: 设$f'(x)$在$x=a$处的右极限存在且有限, 则$f(x)$在$x=a$的右极限也存在且有限.


\begin{proof}

因为$\lim_{x \to a+} f'(x)$存在且有限, 故存在$\delta> 0$, 使得$f'(x)$在$(a, a+\delta)$内有界, 设
\[
|f'(x)| \leq M.
\]
则$\forall \epsilon >0$, 取$\eta = \min \{ \frac{\epsilon}{M}, \delta \}$, 当$x', x'' \in (a, a + \eta)$时,
\[
\Big| \frac{f(x') - f(x'')}{x' - x''} \Big| = |f'(\xi)| \leq M,
\]
$\xi$位于$x', x''$之间. 所以
\[
|f(x') - f(x'')| \leq M|x' - x''| \leq M \cdot \frac{\epsilon}{M} = \epsilon.
\]
由Cauchy收敛准则知$\lim_{x \to a+}f(x)$存在且有限.
\end{proof}



\item 设函数$f(x)$在闭区间$[a, b]$上连续, 在开区间$(a, b)$内可导, 且$f(a) = f(b)$. 求证: 若$f(x)$在$[a, b]$上不为常数, 则$\exists \xi, \eta \in (a, b)$, 使得
\[
f'(\xi) \cdot f'(\eta) <0.
\]


\begin{proof}

由于$f(x)$不为常数, 则$\exists x_0 \in (a, b)$使得$f(x_0) \neq f(a)$. 不妨设
\[
f(x_0) > f(a) = f(b).
\]
分别在$[a, x_0]$和$[x_0, b]$上使用Lagrange中值定理, $\xi \in (a, x_0)$和$\eta \in (x_0, b)$使得
\[
f(x_0) - f(a) = f'(\xi) (x_0 - a), \ f(x_0) - f(b) = f'(\eta) (x_0 - b).
\]
两式相乘, 得
\[
f'(\xi) (x_0 - a) \cdot f'(\eta) (x_0 - b) >0,
\]
注意$(x_0 - a) \cdot (x_0 - b) < 0$, 即证.
\end{proof}
\end{enumerate}

\clearpage

\section{12月21--22日这一周}

\subsection*{题目}

\begin{enumerate}

\item 将$f(x) = \ln \frac{3+x}{2-x}$展开为Maclaurin式(到$x^n$).

\item 将$f(x) = \ln \frac{\sin x}{x}$的Maclaurin式展开到$x^6$项, 这里$f(0)$定义为$x \to 0$时的极限值$f(0) = 0$.

\item 设$f(x)$在$[0, + \infty)$上具有连续的二阶导数, 且
\[
f(0) >0, f'(0) <0, f''(x) < 0 (x \in [0, + \infty))
\]
证明: 存在$\xi \in \left( 0, - \frac{f(0)}{f'(0)}\right)$, 使得$f(\xi ) = 0$.

\item 求常数$a$和$n$, 使得当$x \to 0$时$e^{-x^2} - \cos(\sqrt{2}x) \sim ax^n$.

\item 求极限$\lim_{x \to 0} \frac{x^2 e^{2x} + \ln (1 - x^2)}{x \cos x - \sin x}$.

\item 求极限$\lim_{x \to +\infty} x^{\frac{3}{2}} (\sqrt{x+1} + \sqrt{x-1} - 2 \sqrt x)$.

\item 求函数$f(x) = x^2 2^x$在$x = 0$处的$n$阶导数$f^{(n)}(0)$.

\item 设函数$f(x)$在$[0, 2]$上二阶可导, 且在$[0, 2]$上$|f(x)| \leq a$, $|f''(x)| \leq b$, 证明在$[0, 2]$上成立$|f'(x)| \leq a+b$.

\item 设函数$f(x)$在$[0, 1]$上有三阶导数, 且
\[
f(0) = 0, f(1) = \frac{1}{2}, f'\left( \frac{1}{2} \right) =0.
\]
证明: $\exists \zeta \in (0, 1)$, 使得$f'''(\zeta) \geq 12$.

\end{enumerate}

\clearpage

\subsection*{原稿答案}

\noindent\emph{以下为按题干归入本周的原稿部分参考答案。}

\begin{enumerate}
\item $\lim_{n \to \infty} n[\arctan[\ln(n+1)] - \arctan (\ln n)]$.

\begin{proof}

记$f(x) = \arctan (\ln x)$, 根据Lagrange中值定理
\[
f(n+1) - f(n) = f'(\xi_n) = \frac{1}{(1 + \ln^2 \xi_n) \xi_n}, n < \xi_n < n+1.
\]
\[
\frac{n}{[1 + \ln^2(n+1)](n+1)} < n f'(\xi_n) = \frac{n}{(1 + \ln^2 \xi_n) \xi_n} < \frac{n}{(1 + \ln^2 n) n}.
\]
由此知原极限为0.

\end{proof}


\item 将$f(x) = \ln \frac{3+x}{2-x}$展开为Maclaurin式(到$x^n$).

\begin{proof}
\[
\ln \frac{3+x}{2-x} = \ln \frac{3}{2} + \ln (1 + \frac{x}{3}) - \ln (1 - \frac{x}{2})
\]
\[
=  \ln \frac{3}{2} + \sum_{k=1}^n \frac{1}{k} \left( \frac{1}{2^k} + \frac{(-1)^{k-1}}{3^k} \right)x^k + o(x^n).
\]
\end{proof}


\item 将$f(x) = \ln \frac{\sin x}{x}$的Maclaurin式展开到$x^6$项, 这里$f(0)$定义为$x \to 0$时的极限值$f(0) = 0$.

\begin{proof}

首先利用$\sin x$的展开式
\[
\sin x = x - \frac{1}{3!}x^3 + \frac{1}{5!}x^5 - \frac{1}{7!} x^7 + o (x^8).
\]
就有
\[
\frac{\sin x}{x} = 1 - \frac{1}{3!}x^2 + \frac{1}{5!}x^4 - \frac{1}{7!} x^6 + o (x^7).
\]
因此
\begin{align*}
\ln \frac{\sin x}{x} & = \ln \left(1 - \frac{1}{3!}x^2 + \frac{1}{5!}x^4 - \frac{1}{7!} x^6 + o (x^7) \right)\\
& = \left(- \frac{1}{3!}x^2 + \frac{1}{5!}x^4 - \frac{1}{7!} x^6 + o (x^7) \right)\\
& - \frac{1}{2} \left(- \frac{1}{3!}x^2 + \frac{1}{5!}x^4 - \frac{1}{7!} x^6 + o (x^7) \right)^2 \\
& + \frac{1}{3} \left(- \frac{1}{3!}x^2 + \frac{1}{5!}x^4 - \frac{1}{7!} x^6 + o (x^7) \right)^3 + o(x^7)\\
& =  \left(- \frac{1}{3!}x^2 + \frac{1}{5!}x^4 - \frac{1}{7!} x^6 + o (x^7) \right) - \frac{1}{2} \left(\frac{1}{36} x^4 - \frac{2}{3! 5!}x^6 \right) + \frac{1}{3} \left( - \frac{1}{216} x^6 \right) + o(x^7) \\
& = -\frac{1}{6}x^2 - \frac{1}{180}x^4 - \frac{1}{2835}x^6 + o(x^7).
\end{align*}

\end{proof}


\item 设$f(x)$在$[0, + \infty)$上具有连续的二阶导数, 且
\[
f(0) >0, f'(0) <0, f''(x) < 0 (x \in [0, + \infty))
\]
证明: 存在$\xi \in \left( 0, - \frac{f(0)}{f'(0)}\right)$, 使得$f(\xi ) = 0$.

\begin{proof}
$\forall x \in [0, + \infty)$, 在$0$点展开为Lagrange余项:
\[
f(x) = f(0) + f'(0)x + \frac{1}{2} f''(\theta x) x^2, \theta \in (0, 1).
\]
两边代入$x = - \frac{f(0)}{f'(0)}$, 并利用条件
\[
f\left(- \frac{f(0)}{f'(0)} \right) = f(0) - f(0) +  \frac{1}{2} f''(\theta x) \left( - \frac{f(0)}{f'(0)}\right)^2 <0,
\]
利用零点定理即得.
\end{proof}


\item 求函数$f(x) = x^2 2^x$在$x = 0$处的$n$阶导数$f^{(n)}(0)$.


\begin{proof}

可以用Leibniz定理, 也可以用Taylor展开. 由于
\[
2^x = e^{x \ln 2} = 1 + x \ln 2 + \frac{(x \ln 2)^2}{2} + \cdots + \frac{1}{(n-2)!} (x \ln 2)^{n-2} + o(x^{n-2}).
\]
故
\[
f(x) = x^2 2^x = x^2 + x^3 \ln 2 + \frac{x^4 (\ln 2)^2}{2} + \cdots + \frac{1}{(n-2)!} x^n (\ln 2)^{n-2} + o(x^n).
\]
因此
\[
\frac{f^{(n)}(0)}{n!} = \frac{1}{(n-2)!} (\ln 2)^{n-2}, n \geq 2
\]
可求得
\[
f^{(n)}(0) = n (n-1)(\ln 2)^{n-2}, n \geq 2.
\]

\end{proof}

\item 设函数$f(x)$在$[0, 2]$上二阶可导, 且在$[0, 2]$上$|f(x)| \leq a$, $|f''(x)| \leq b$, 证明在$[0, 2]$上成立$|f'(x)| \leq a+b$.


\begin{proof}

对$\forall x_0 \in [0, 2]$, 对于$x =0$和$x=2$分别展开为Lagrange余项的Taylor公式:
\[
f(0) = f(x_0) + f'(x_0) (0 - x_0) + \frac{f''(\zeta_1)}{2} (0 - x_0)^2, \zeta_1 \in (0, x_0),
\]
\[
f(2) = f(x_0) + f'(x_0) (2 - x_0) + \frac{f''(\zeta_2)}{2} (2 - x_0)^2, \zeta_2 \in (x_0, 2).
\]
两式作差取绝对值, 可得
\[
|f'(x_0)| = \frac{1}{2} \Big| f(2) - f(0) + \frac{f''(\zeta_1)}{2}x_0^2 - \frac{f''(\zeta_2)}{2} (x_0 - 2)^2 \Big|
\]
\[
\leq \frac{1}{2} \left( |f(0)| + |f(2)| + \frac{|f''(\zeta_1)|}{2}x_0^2 +  \frac{|f''(\zeta_2)|}{2} (x_0 - 2)^2  \right)
\]
\[
\leq a + \frac{b}{4} (x_0^2 + (x_0 - 2)^2) \leq a+b.
\]
注意函数$x_0^2 + (x_0 - 2)^2$在$[0, 2]$上的最大值为4.
\end{proof}



\item 设函数$f(x)$在$[0, 1]$上有三阶导数, 且
\[
f(0) = 0, f(1) = \frac{1}{2}, f'\left( \frac{1}{2} \right) =0.
\]
证明: $\exists \zeta \in (0, 1)$, 使得$f'''(\zeta) \geq 12$.


\begin{proof}

展开为Lagrange余项的Taylor公式:
\[
f(1) = f(\frac{1}{2}) + f'(\frac{1}{2}) (1 - \frac{1}{2}) + \frac{1}{2} f''(\frac{1}{2}) (1 - \frac{1}{2})^2 + \frac{1}{6} f'''(\xi) (1 - \frac{1}{2})^3, \xi \in (\frac{1}{2}, 1),
\]
\[
f(0) = f(\frac{1}{2}) +f'(\frac{1}{2}) (0 - \frac{1}{2}) + \frac{1}{2} f''(\frac{1}{2}) (0 - \frac{1}{2})^2 + \frac{1}{6} f'''(\eta) (0 - \frac{1}{2})^3, \eta \in (0, \frac{1}{2}).
\]
两式想减并化简, 得到
\[
f'''(\xi) + f'''(\eta) = 24.
\]
因此$f'''(\xi)$和$f'''(\eta)$中至少一个$\geq 12$.

\end{proof}


\end{enumerate}

\clearpage

\section{12月28--29日这一周}

\subsection*{题目}

\begin{enumerate}

\item 设$n \in \mathbb N^*$, 证明:
\[
e^x > 1 + \frac{x}{1!} + \frac{x^2}{2!} + \cdots + \frac{x^n}{n!}
\]
对$x > 0$成立.

\item 求数列$\{\sqrt[n]{n} \}$中的最大项.

\item 设$x \in (0, \frac{\pi}{2})$, 证明: $(\sin x)^{1 - \cos 2x} + (\cos x)^{1 + \cos 2x} \geq \sqrt 2$.

\item 设$f(x)$在$(a, b)$上为凸函数, 且有界. 证明: $f(a+)$和$f(b-)$存在.

\item 已知$a > 0$, 求函数$y = (1 - x)e^x$在闭区间$[-a, a]$上的最值.

\item 证明$4 \arctan x - x + \frac{4 \pi}{3} - \sqrt 3 =0$恰有两个实根.

\item 设$f(x)$在$(-\infty, +\infty)$内可导, 且$\phi(x) = \frac{f(x)}{x}$在$x = a (a \neq 0)$处取得极值. 证明: 曲线$y = f(x)$在点$(a, f(a))$处的切线经过原点.

\end{enumerate}

\clearpage

\subsection*{原稿答案}

\noindent\emph{原稿未附这一周对应的参考答案。}

\clearpage

\chapter{春季学期}

\clearpage

\section{3月25--31日 · Week 5}

\subsection*{题目}

\begin{definition}
设$f(x)$为$[a, b]$上的有限函数, 如果对$[a, b]$的一切分割
\[
\pi: a = x_0  < x_1 < x_2 < \cdots < x_{n-1} < x_n =b,
\]
使集合$\{ \sum_{i=1}^n |f(x_i) - f(x_{i-1})| \}$为有界集, 即
\[
M = \sup_{\pi}\{ \sum_{i=1}^n |f(x_i) - f(x_{i-1})| \} < \infty,
\]
则称$f(x)$是$[a, b]$上的有界变差函数, 并称$M$为$f(x)$在$[a, b]$上的全变差, 记为$V_a^b(f)$.
\end{definition}

\begin{enumerate}
\item 证明$[a, b]$上的有界变差函数是有界的.
\item 证明$[a, b]$上的单调函数是有界变差的.
\item ({\bf Jordan分解定理}) 证明$f$是$[a, b]$上的有界变差函数的充分必要条件是$f$可以分解为两个单调增函数的差. (tip: 证明函数$V_a^x(f)$是单调的).
\item 证明$[a, b]$上的有界变差函数Riemann可积.
\end{enumerate}

\begin{theorem}[加强形式的有限覆盖定理]
证明: 若$\{ \mathcal O_{\alpha} \}$是$[a, b]$的一族开覆盖, 则存在一个正数$\delta >0$, 使得对于区间$[a, b]$中的任何两个点$x', x''$, 只要$|x' - x''|< \delta$, 就存在开覆盖族中的一个开区间, 覆盖$x', x''$. (数$\delta$称为该开覆盖的Lebesgue数).
\end{theorem}

\begin{enumerate}

\item 证明上述定理.

\item 设函数$f$在$[a, b]$上有界, 如果$f$的所有不连续点可以用长度任意小的至多可数个开区间覆盖, 则$f$可积.

\item 若函数$f(x), g(x)$在$[0, 1]$上有相同的单调性. 证明:
\[
\int_0^1 f(x) dx \int_0^1 g(x) dx \leq \int_0^1 f(x) g(x) dx.
\]
(tip: 利用排序不等式).

\item 按照以下步骤证明$\pi$是无理数:

\begin{itemize}
\item[Step 1.] 反证, 设$\pi = \frac{a}{b}, (a, b)=1$.
\item[Step 2.] 作函数
\[
f(x) = \frac{1}{n!} x^n (a-bx)^n, \ g(x) = f(x) - f^{(2)}(x) + f^{(4)}(x) - \cdots + (-1)^n f^{(2n)}(x).
\]
\item[Step 3.] 证明: $[g'(x) \sin x - g(x) \cos x]' = f(x) \sin x$.
\item[Step 4.] 先说明$f(\pi-x) = f(x)$, 再证明$g(\pi) = g(0)$是整数.
\item[Step 5.] 让$n \to \infty$, 尝试得到矛盾.
\end{itemize}

\end{enumerate}

\clearpage

\subsection*{原稿答案}

\noindent\emph{以下为按题干归入本周的原稿部分参考答案。}

\begin{enumerate}
\item 证明$[a, b]$上的有界变差函数是有界的.

\begin{proof}

对于任意$a \leq x \leq b$, 得到分割$[a, x]$和$[x, b]$. 根据有界变差的定义,
\[
|f(x) - f(a)| + |f(b) - f(x)| \leq M.
\]
显然
\[
|f(x)| \leq |f(a)| + |f(a) - f(x)| \leq |f(a)| + M.
\]

\end{proof}


\item 证明$[a, b]$上的单调函数是有界变差的.

\begin{proof}

不妨设函数单调增. 任取分割$\pi: a = x_0 < x_1 < \cdots < x_n = b$, 有
\[
\sum_{i=1}^n |f(x_i) -f(x_{i-1})| = f(b) - f(a).
\]
即证.

\end{proof}

\item ({\bf Jordan分解定理}) 证明$f$是$[a, b]$上的有界变差函数的充分必要条件是$f$可以分解为两个单调增函数的差. (tip: 证明函数$V_a^x(f)$是单调的).

\begin{proof}
\

\noindent
{\bf 必要性:} 分解$f(x) = V_a^x (f) - [V_a^x (f) - f(x)]$. 易证$V_a^x (f)$为单调增函数. 再证明$g(x) = V_a^x (f) - f(x)$也是单调增函数. 设$x_1 < x_2$, 
\[
g(x_2) - g(x_1) = | V_a^{x_2}(f) -( f(x_2) - f(a)) | - | V_a^{x_1}(f) - (f(x_1) - f(a)) | = V_{x_1}^{x_2}(f) - [f(x_2) - f(x_1)] \geq 0.
\]



\noindent
{\bf 充分性:} 利用上一问, 只需证明有界变差函数之差仍为有界变差函数.


\end{proof}


\item 证明$[a, b]$上的有界变差函数Riemann可积.


\begin{proof}
利用上一问显然.
\end{proof}



\end{enumerate}


\begin{enumerate}


\item 证明加强形式的有界覆盖定理.


\begin{proof}

首先用覆盖定理得到区间$[a, b]$的一个有限覆盖$I_1, I_2, \cdots, I_k$. 将这有限个开区间所有端点从小到大排列, 去掉重复点, 记为
\[
a = x_0 < x_1 < x_2 < \cdots < x_N =b.
\]
记这个端点集合为$A = \{x_0, x_1, \cdots, x_N \}$. 令
\[
\delta = \min \{ x_1 - x_0, x_2 - x_1, \cdots, x_N - x_{N-1} \}.
\]
再证明该$\delta$即为所需的Lebesgue数.


任取两点$x', x'' \in [a, b]$, 使得$0 < x'' - x ' < \delta$. 则分两种情况:

(1) 在闭区间$[x', x'']$内没有集合$A$中的点, 覆盖该闭区间中任何一点的开区间也覆盖了整个闭区间.


(2) 在闭区间$[x', x'']$内有$A$中的点. 由于$\delta$的取法, 这样的点只有一个, 记为$x_i$. 则$x' \leq x_i \leq x''$. 那么覆盖$x_i$的开区间(并非以$x_i$为端点的小区间)一定覆盖$x'$和$x''$.


\end{proof}

\item 设函数$f$在$[a, b]$上有界, 如果$f$的所有不连续点可以用长度任意小的至多可数个开区间覆盖, 则$f$可积.


\begin{proof}

对于任给的$\epsilon, \eta > 0$, 先用总长度小于$\eta$的有限个或者可数个开区间覆盖$f$的所有不连续点. 根据连续性定义, 对于每个连续点$x_0$, 都可以找到一个邻域$O(x_0)$, 使得$f$在这个邻域上振幅小于$\epsilon$. 所有这些连续点的邻域和覆盖不连续点的开区间一起构成区间$[a, b]$的一个开覆盖$\{\mathcal O_{\alpha} \}$.


利用加强形式的覆盖定理, 存在Lebesgue数$\delta >0$, 使得在区间$[a, b]$中相距不超过$\delta$的任何两个点可以用开覆盖$\{\mathcal O_{\alpha} \}$中的某一个开区间同时覆盖.


对于区间$[a, b]$取细度不超过$\delta$的等距分割$\pi$. 这时每个子区间被开覆盖中的一个开区间所覆盖. 如果子区间是由原来覆盖不连续点的开区间所覆盖, 则这类子区间的总长度必小于$\eta$. 而覆盖所有其他子区间的开区间都是原先用于覆盖连续点的邻域, 因此$f$在每一个这类子区间上的振幅小于$\epsilon$.


最后应用第三充要条件即证.


\end{proof}



\item 若函数$f(x), g(x)$在$[0, 1]$上有相同的单调性. 证明:
\[
\int_0^1 f(x) dx \int_0^1 g(x) dx \leq \int_0^1 f(x) g(x) dx.
\]

\begin{proof}

因为$f(x), g(x)$在$[0, 1]$上可积, 故$f(x) \cdot g(x)$在$[0, 1]$上可积. 在$[0, 1]$上作等距分割:
\[
\pi: 0 = x_0 < x_1 < x_2 < \cdots < x_n =1,
\]
使得$\Delta x_k = \left[ \frac{k-1}{n}, \frac{k}{n} \right]$. 取$\xi_k = \frac{k}{n}, 1 \leq k \leq n$. 因为$f(x), g(x)$在$[0, 1]$上有相同的单调性, 故
\[
\sum_{k=1}^n f \left( \frac{k}{n} \right) \sum_{k=1}^n g \left( \frac{k}{n} \right) \leq n \sum_{k=1}^n f \left( \frac{k}{n} \right) g \left( \frac{k}{n} \right).
\]
上式两边同除以$n^2$, 即得
\[
\frac{1}{n} \sum_{k=1}^n f \left( \frac{k}{n} \right) \cdot \frac{1}{n} \sum_{k=1}^n g \left( \frac{k}{n} \right) \leq \frac{1}{n} \sum_{k=1}^n f \left( \frac{k}{n} \right) g \left( \frac{k}{n} \right).
\]
根据定积分的定义, 令$n \to \infty$取极限即证.
\end{proof}


此处将左边求和式乘积展开, 再分成$n$组, 之后利用排序不等式:


\begin{formula}[排序不等式]
设有两列数列$a_1 \leq a_2 \leq \cdots \leq a_n$和$b_1 \leq b_2 \leq \cdots \leq b_n$, 则
\[ a_1 b_1 + a_2 b_2 + \cdots + a_n b_n \geq a_1 b_{\sigma(1)} + a_2 b_{\sigma(2)} + \cdots + a_n b_{\sigma(n)} \geq a_1 b_n + a_2 b_{n-1} + \cdots + a_n b_1,
\]
其中$\sigma(1), \sigma(2), \ldots, \sigma(n)$是$1,2, \ldots, n$的任意置换(即顺序和$\geq$乱序和$\geq$逆序和). 其中等号成立当且仅当$a_1 =a_2 = \cdots = a_n$或$b_1 =b_2 = \cdots = b_n$.
\end{formula}



\item 按照步骤证明$\pi$是无理数.


\begin{proof}
反证, 设$\pi$是有理数: $\pi = \frac{a}{b}, (a, b)=1$. 作$f(x) = \frac{1}{n!}x^n (a - bx)^n$,
\[
g(x) = f(x) - f^{(2)}(x) + f^{(4)}(x) - \cdots + (-1)^n f^{(2n)}(x).
\]
显然, $f^{(2n+1)}(x) \equiv 0$, 且
\[
g''(x) = f^{(2)}(x) - f^{(4)}(x) + \cdots.
\]
故
\[
g''(x) + g(x) = f(x) \mbox{ 及 \ } [g'(x) \sin x - g(x) \cos x]' = [g''(x) + g(x)] \sin x = f(x) \sin x. 
\]
在上式两边作$[0, \pi]$上积分, 则有$\int_0^{\pi} f(x) \sin x dx = g(\pi) + g(0)$. 由于
\[
f(\pi - x) = f \left( \frac{a}{b} - x \right) = \frac{1}{n!} \left( \frac{a}{b} - x \right)^n \left[a - b \left( \frac{a}{b} - x \right) \right]^n = f(x),
\]
所以得到
\[
f^{(2k)} (\pi - x) = f^{(2k)} (x), f(\pi) = f(0), f^{(2k)}(\pi) = f^{(2k)}(0).
\]
从而推出$g(\pi) = g(0)$是整数, 即$\int_0^{\pi} f(x) \sin x dx$是整数.


由$0 < \sin x \leq 1 (0 < x < \pi)$, $0 < \frac{1}{n!}x^n (a - bx)^n < \frac{\pi^n a^n}{n!}$, 知$0 < f(x) \sin x < \frac{\pi^n a^n}{n!}$. 从而得到
\[
0 < \int_0^{\pi} f(x) \sin x dx < \frac{\pi^n a^n}{n!} \pi.
\]
矛盾.

\end{proof}



\end{enumerate}

\clearpage

\section{4月1--7日 · Week 6}

\subsection*{题目}

\begin{enumerate}

\item 证明Minkowski不等式:
设$f(x), g(x)$在$[a, b]$上连续, $1 \leq p < +\infty$, 则
\[
\left( \int_a^b |f(x) + g(x)|^p dx \right)^{\frac{1}{p}} \leq \left( \int_a^b |f(x)|^p dx \right)^{\frac{1}{p}} + \left( \int_a^b |g(x)|^p dx \right)^{\frac{1}{p}}.
\]
当$0 < p < 1$时不等式反向成立.

\item 设$f(x)$在$[0, 1]$上有连续导数, 且$f(1) - f(0) =1$. 求证: 
\[
\int_0^1 (f'(x))^2 dx \geq 1.
\]

\item 设$f(x)$是以$T$为周期的连续函数, $\exists x_0$, 使得$f(x_0) \neq 0$, 且$\int_0^T f(x) dx =0$. 求证: 在$(x_0, x_0 + T)$内至少存在两个零点.

\item 设$f(x) \geq 0, g(x) \geq 0$, 且在$[a, b]$上连续. 求证: 
\[
\lim_{n \to \infty} \sqrt[n]{\int_a^b (f(x))^n g(x) dx } = \max_{a \leq x \leq b} f(x).
\]

\item 证明Riemann引理:
设$f \in R[a, b]$, 则
\[
\lim_{p \to +\infty} \int_a^b f(x) \sin px dx = 0, \ \ \lim_{p \to +\infty} \int_a^b f(x) \cos px dx = 0.
\]

\item 设函数$f$在闭区间$[A, B]$上可积. 证明$f$具有积分的连续性, 即
\[
\lim_{h \to 0} \int_a^b | f(x+h) - f(x) | dx =0 \ (A < a < b < B).
\]

Tip: 证法一: 利用已知结论, 若$f$在$[a, b]$上可积, 则$\forall \epsilon >0$, 存在连续函数$g(x)$使得$\int_a^b |f(x) - g(x)| dx < \epsilon$. 证法二: 对区间作等距分割, 转化为振幅的讨论.

\end{enumerate}

\clearpage

\subsection*{原稿答案}

\noindent\emph{以下为按题干归入本周的原稿部分参考答案。}

\begin{enumerate}

\item 证明Minkowski不等式:
设$f(x), g(x)$在$[a, b]$上连续, $1 \leq p < +\infty$, 则
\[
\left( \int_a^b |f(x) + g(x)|^p dx \right)^{\frac{1}{p}} \leq \left( \int_a^b |f(x)|^p dx \right)^{\frac{1}{p}} + \left( \int_a^b |g(x)|^p dx \right)^{\frac{1}{p}}.
\]
当$0 < p < 1$时不等式反向成立.


\begin{proof}
利用H\textrm{$\ddot{o}$}lder不等式:
\[
\int_a^b |f(x) + g(x)|^p dx = \int_a^b |f(x) + g(x)| |f(x) + g(x)|^{p-1} dx
\]
\[
\leq \int_a^b |f(x)| |f(x) + g(x)|^{p-1} dx + \int_a^b | g(x)| |f(x) + g(x)|^{p-1} dx
\]
\[
\leq \left( \int_a^b |f(x)|^p dx \right)^{\frac{1}{p}} \left( \int_a^b |f(x) + g(x)|^{q(p-1)} dx\right)^{\frac{1}{q}} + \left( \int_a^b |g(x)|^p dx \right)^{\frac{1}{p}} \left( \int_a^b |f(x) + g(x)|^{q(p-1)} dx\right)^{\frac{1}{q}} 
\]
\[
= \left[ \left( \int_a^b |f(x)|^p dx \right)^{\frac{1}{p}} +  \left( \int_a^b |g(x)|^p dx \right)^{\frac{1}{p}} \right]\left( \int_a^b |f(x) + g(x)|^{q(p-1)} dx\right)^{\frac{1}{q}}.
\]
注意$p = q(p-1)$且$1 - \frac{1}{q} = \frac{1}{p}$, 即证.
\end{proof}


\item 设$f(x)$在$[0, 1]$上有连续导数, 且$f(1) - f(0) =1$. 求证: 
\[
\int_0^1 (f'(x))^2 dx \geq 1.
\]


\begin{proof}

利用Cauchy-Schwarz公式即证. 注意$\int_0^1 f'(x)d x = f(1) - f(0) = 1$.

\end{proof}


\item 设$f(x)$是以$T$为周期的连续函数, $\exists x_0$, 使得$f(x_0) \neq 0$, 且$\int_0^T f(x) dx =0$. 求证: 在$(x_0, x_0 + T)$内至少存在两个零点.


\begin{proof}
若无零点, 则$f(x)$恒正或恒负, 因此得到矛盾
\[
\int_0^T f(x) dx  = \int_{x_0}^{x_0 +T} f(x) dx \neq 0.
\]
故必$\exists \xi \in (x_0, x_0 + T)$使得$f(\xi) =0$. 根据周期性, $f(\xi +T)=0$. 但注意$\xi + T \notin (x_0, x_0 + T)$.


在区间$(\xi, \xi+T)$上重复上述讨论, 存在零点$\xi_1 \in (\xi, \xi+T)$. 易证$\xi_1$或$\xi_1-T$必有其一在$(x_0, x_0 + T)$, 且为零点. 即证.

\end{proof}



\item 设$f(x) \geq 0, g(x) \geq 0$, 且在$[a, b]$上连续. 求证: 
\[
\lim_{n \to \infty} \sqrt[n]{\int_a^b (f(x))^n g(x) dx } = \max_{a \leq x \leq b} f(x).
\]

\begin{proof}

根据连续性, $\exists x_0 \in [a, b]$, 使得$f(x_0) =  \max_{a \leq x \leq b} f(x) : = M$. 则
\[
\left[ \int_a^b (f(x))^n g(x) dx \right]^{\frac{1}{n}} \leq \left( \int_a^b M^n g(x) dx \right)^{\frac{1}{n}} = M \left( \int_a^b g(x) dx \right)^{\frac{1}{n}}.
\]
两端取极限, 得
\[
\lim_{n \to \infty} \sqrt[n]{\int_a^b (f(x))^n g(x) dx } \leq M.
\]
另一方面, 因为$f(x)$在$x_0$连续, 所以$\forall \epsilon >0, \exists [\alpha, \beta] \subset [a, b]$,使得
\[
f(x) > M - \epsilon, \ \ \forall x \in [\alpha, \beta].
\]
从而
\[
\left[ \int_a^b (f(x))^n g(x) dx \right]^{\frac{1}{n}}  \geq \left[ \int_{\alpha}^{\beta} (f(x))^n g(x) dx \right]^{\frac{1}{n}} \geq (M - \epsilon ) \left( \int_{\alpha}^{\beta} g(x) dx \right)^{\frac{1}{n}}.
\]
两端令$n \to \infty$, 得
\[
\lim_{n \to \infty} \sqrt[n]{\int_a^b (f(x))^n g(x) dx } \geq M- \epsilon.
\]
综合即证.

\end{proof}


\item 证明Riemann引理:
设$f \in R[a, b]$, 则
\[
\lim_{p \to +\infty} \int_a^b f(x) \sin px dx = 0, \ \ \lim_{p \to +\infty} \int_a^b f(x) \cos px dx = 0.
\]

\begin{proof}


对任意有界区间$[\alpha, \beta]$有
\[
\Big| \int_{\alpha}^{\beta} \sin px dx \Big| = \Big| \frac{\cos p \alpha - \cos p \beta}{p} \Big| \leq \frac{2}{p}.
\]
设在$[a, b]$上$|f(x) | \leq M$. 任给$\epsilon >0$, 则存在$[a, b]$的分割$\pi$:
\[
a = x_0 < x_1 < \cdots < x_n =b,
\]
使得$\sum \omega_i \Delta x_i < \frac{\epsilon}{2}$. 于是当$p \geq 4n \frac{M}{\epsilon}$时, 有
\[
\Big| \int_a^b f(x) \sin px dx \Big| = \Big| \sum_{k=1}^n \int_{x_{k-1}}^{x_k} [f(x_k)  + f(x) - f(x_k)]\sin px dx \Big|
\]
\[
\leq \sum_{k=1}^n \{ |f(x_k)| \Big| \int_{x_{k-1}}^{x_k} \sin px dx \Big| + \int_{x_{k-1}}^{x_k} |f(x) - f(x_k)| |\sin px| dx \}
\]
\[
< \frac{2nM}{p} +  \sum \omega_i \Delta x_i  < \epsilon.
\]

\end{proof}


\item 设函数$f$在闭区间$[A, B]$上可积. 证明$f$具有积分的连续性, 即
\[
\lim_{h \to 0} \int_a^b | f(x+h) - f(x) | dx =0 \ (A < a < b < B).
\]

\begin{proof}
\

\noindent
{\bf 证一:} 对任给的$\epsilon > 0$, 因$f$在$[A, B]$上可积, 故存在$[A, B]$上的连续函数$\varphi$, 使得
\[
\int_A^B |f(x) - \varphi(x)|dx < \frac{\epsilon}{4} \ (A < a < b < B).
\]
由于$\varphi$在$[A, B]$上一致连续, 故存在$\delta >0$, 使当$x', x'' \in [A, B], |x' - x''| < \delta$时, 恒有
\[
|\varphi(x') - \varphi(x'')| < \frac{\epsilon}{2(b-a)}.
\]
于是当$|h| < \delta$时,
\[
\int_a^b | f(x+h) - f(x) | dx  \leq \int_a^b | f(x+h) - \varphi(x+h) | dx + \int_a^b | \varphi(x+h) - \varphi(x)|dx + \int_a^b |\varphi(x) - f(x)|dx
\]
\[
\leq 2 \int_A^B |f(x) - \varphi(x)|dx + \int_a^b | \varphi(x+h) - \varphi(x) | dx  < 2 \cdot \frac{\epsilon}{4} + \frac{\epsilon}{2(b-a)} (b-a) = \epsilon.
\]


\noindent
{\bf 证二:} 作$[a, b]$的等距分割$\pi: x_i = a_i + i\frac{b-a}{n}$. 则
\[
\int_a^b | f(x+h) - f(x) | dx = \sum_{i=1}^n \int_{x_{i-1}}^{x_i} |f(x+h) - f(x) | dx.
\]
若令$|h| < \frac{b-a}{n}$, 则点$x+h$要么位于$x$所在的第$i$个小区间上, 从而
\[
 | f(x+h) - f(x) | < \omega_i.
\]
要么点$x+h$在相邻的一个区间上, 此时
\[
 | f(x+h) - f(x) |  < \omega_i + \omega_{i+1} (h >0) \mbox{ 或 \ } \omega_{i-1} + \omega_i (h < 0).
\]
总之有$|f(x+h) - f(x)| \leq \omega_{i-1} + \omega_i + \omega_{i+1}$. 对第$1, n$两个小区间只要将$\omega_0, \omega_{n+1}$看作$2M$即可. 于是
\[
\sum_{i=1}^n \int_{x_{i-1}}^{x_i} |f(x+h) - f(x) | dx \leq 3 \sum_{i=1}^n \omega_i \Delta_i + 2M \frac{b-a}{n}.
\]
根据可积性, $\forall \epsilon>0$, 只要$n$充分大, 即可使得
\[
\sum_{i=1}^n \omega_i \Delta_i  < \frac{\epsilon}{6}, \ \ 2M \frac{b-a}{n} < \frac{\epsilon}{2}.
\]
取$|h| < \delta < \frac{b-a}{n}$即可.

\end{proof}

\begin{proposition}[Riemann-Lebesgue定理]
设$f(x)$在区间$[a, b]$上具有一阶连续导数, 证明:
\[
\lim_{p \to + \infty} \int_a^b f(x) \sin(px) dx = 0.
\]
\end{proposition}


\begin{proof}
利用分部积分得
\[
\int_a^b f(x) \sin(px) dx = -\frac{1}{p} (f(b) \cos(pb) - f(a) \cos(pa)) + \frac{1}{p} \int_a^b f'(x) \cos (px) dx,
\]
所以
\[
\Big| \int_a^b f(x) \sin(px) dx \Big| \leq \frac{1}{|p|}[|f(b)| + |f(a)|] + \frac{1}{|p|} \int_a^b |f'(x)|dx,
\]
令$p \to + \infty$即证. 类似证明
\[
\lim_{p \to + \infty} \int_a^b f(x) \cos(px) dx = 0.
\]

\end{proof}



\end{enumerate}

\clearpage

\section{4月22--28日 · Week 9}

\subsection*{题目}

\begin{enumerate}
\item 求极限
\[
\lim_{n \to \infty} \sum_{k = 1}^n \left( k \ln \frac{2k+1}{2k-1} -1 \right).
\]

\begin{definition}[幂平均]
设$a_i \geq 0 \ (1 \leq i \leq n)$, 记
\[
M(r) = \left( \frac{1}{n} \sum_{i=1}^n a_i^r \right)^{\frac{1}{r}} \ (r > 0),
\]
称$M(r)$为$a_1, \cdots, a_n$的$r$次幂平均. 它与算术平均(AM)的关系为
\[
M(1) = \frac{a_1 + a_2 + \cdots + a_n}{n} = AM,
\]
\[
M(r) =( AM(a_1^r, a_2^r, \cdots, a_n^r))^{\frac{1}{r}}.
\]
\end{definition}

\begin{definition}[加权平均]
设$p_i > 0 \ (1 \leq i \leq n)$, 记
\[
M_p(r) = \left( \frac{ \sum_{i=1}^n p_i a_i^r }{\sum_{i=1}^n p_i} \right)^{\frac{1}{r}}, 
\]
\[
G_p = \left( \prod_{i=1}^n a_i^{p_i} \right)^{1 / \sum_{i=1}^n p_i} = (a_1^{p_1} a_2^{p_2} \cdots a_n^{p_n} )^{\frac{1}{p_1 + p_2 + \cdots + p_n}}.
\]
$M_p(r)$和$G_p$分别称为$a_1, a_2, \cdots, a_n$的加权($r$次幂)算术平均和加权几何平均, 其中$p_i$称为权数. 若令$q_i = \frac{p_i}{p_1 + \cdots + p_n}$, 则$\sum_{i=1}^n q_i =1$. 这时
\[
M_q(r) = \left( \sum_{i=1}^n q_i a_i^r \right)^{\frac{1}{r}}
\]
\[
G_q = \prod_{i=1}^n a_i^{q_i} = a_1^{q_1} a_2^{q_2} \cdots a_n^{q_n}. 
\]
将$p_i$转化成$q_i$的过程称为权数的标准化.

\end{definition}

\item 设$r > 0, a_1, a_2, \cdots, a_n$不全相等, 证明:
\[
M_q(r) > M_q\left(\frac{r}{2}\right).
\]

\item 证明:
\[
G_q = \lim_{r \to 0+} M_q(r).
\]

\item 设$a_1, a_2, \cdots, a_n$不全相等, 则$G_q < M_q(1)$. 即:
\[
a_1^{q_1} a_2^{q_2} \cdots a_n^{q_n} < q_1 a_1 + \cdots + q_n a_n.
\]
等号成立的条件为$a_1 = a_2 = \cdots = a_n$.

\item 已知椭圆$\frac{x^2}{a^2} + \frac{y^2}{b^2} =1 (a, b >0)$的周长:
\[
s = 4 \int_0^{\frac{\pi}{2}} \sqrt{a^2 \sin^2 t + b^2 \cos^2 t} dt,
\]
证明: $\pi(a+b) \leq s \leq \pi \sqrt{2a^2 + 2b^2}$. Tip: 利用Cauchy-Schwarz不等式积分/离散形式.

\item 设$f(x)$在区间$[a, b]$上可积. 对区间$[a, b]$作等距分割: 
\[
\pi: a = x_0 < x_1 < \cdots < x_n = b, \ \Delta x_i = \frac{b-a}{n}. 
\]
推导此情形下的抛物线法公式.

\item 设$f \in C^2[0, h]$, 证明: 存在$\xi \in [0, h]$使得
\[
\int_0^h f(x) dx = \frac{h}{2} [f(0) + f(h)] - \frac{1}{12} f''(\xi) h^3. 
\]

\item 设$f \in C^4[0, h]$, 证明: 存在$\xi \in [0, h]$使得
\[
\int_0^h f(x) dx = \frac{h}{2} [f(0) + f(h)] - \frac{h^2}{12} [f'(h) - f'(0)] + \frac{1}{720} f^{(4)} (\xi) h^5.
\]

\item 设$f(x)$在$[0, 1]$上可导, 且当$x \in (0, 1)$时, $0 < f'(x) < 1, f(0) = 0$. 试证:
\[
\left( \int_0^1 f(x) dx \right)^2 > \int_0^1 f^3(x) dx.
\]

\end{enumerate}

\clearpage

\subsection*{原稿答案}

\noindent\emph{以下为按题干归入本周的原稿部分参考答案。}

\begin{enumerate}

\item 求极限
\[
\lim_{n \to \infty} \sum_{k = 1}^n \left( k \ln \frac{2k+1}{2k-1} -1 \right).
\]


\begin{proof}


记$S_n =  \sum\limits_{k = 1}^n \left( k \ln \frac{2k+1}{2k-1} -1 \right)$, 则
\[
S_n = n \ln (2n + 1) - \ln [ (2n-1)!! ] - n = n \ln (2n+1) - n - \ln \frac{(2n)!}{2^n \cdot n!}
\]
\[
= n \ln (2n+1) - n - \ln [ (2n)! ] + n \ln 2 + \ln (n!).
\]
利用Stirling公式:
\[
\ln (n!) = n \ln n - n + \frac{\ln n}{2} + \frac{\ln 2}{2} + \frac{\ln \pi}{2} + o(1) \ (n \to \infty),
\]
\[
\ln [ (2n)! ] = 2n \ln 2n - 2n + \frac{\ln 2n}{2} + \frac{\ln 2}{2} +  \frac{\ln \pi}{2} + o(1) \ (n \to \infty),
\]
\[
\ln (2n+1) = \ln 2n + \frac{1}{2n} + o \left( \frac{1}{n} \right) \ (n \to \infty).
\]
从而有$S_n = \frac{1}{2} - \frac{\ln 2}{2} + o(1) \ (n \to \infty)$.
\end{proof}



\begin{definition}[幂平均]
设$a_i \geq 0 \ (1 \leq i \leq n)$, 记
\[
M(r) = \left( \frac{1}{n} \sum_{i=1}^n a_i^r \right)^{\frac{1}{r}} \ (r > 0),
\]
称$M(r)$为$a_1, \cdots, a_n$的$r$次幂平均. 它与算术平均(AM)的关系为
\[
M(1) = \frac{a_1 + a_2 + \cdots + a_n}{n} = AM, \ \ M(r) =( AM(a_1^r, a_2^r, \cdots, a_n^r))^{\frac{1}{r}}.
\]
\end{definition}


\begin{definition}[加权平均]
设$p_i > 0 \ (1 \leq i \leq n)$, 记
\[
M_p(r) = \left( \frac{ \sum_{i=1}^n p_i a_i^r }{\sum_{i=1}^n p_i} \right)^{\frac{1}{r}}, \ \
G_p = \left( \prod_{i=1}^n a_i^{p_i} \right)^{1 / \sum_{i=1}^n p_i} = (a_1^{p_1} a_2^{p_2} \cdots a_n^{p_n} )^{\frac{1}{p_1 + p_2 + \cdots + p_n}}.
\]
$M_p(r)$和$G_p$分别称为$a_1, a_2, \cdots, a_n$的加权($r$次幂)算术平均和加权几何平均, 其中$p_i$称为权数. 若令$q_i = \frac{p_i}{p_1 + \cdots + p_n}$, 则$\sum_{i=1}^n q_i =1$. 这时
\[
M_q(r) = \left( \sum_{i=1}^n q_i a_i^r \right)^{\frac{1}{r}}
\]
\[
G_q = \prod_{i=1}^n a_i^{q_i} = a_1^{q_1} a_2^{q_2} \cdots a_n^{q_n}. 
\]
将$p_i$转化成$q_i$的过程称为权数的标准化.

\end{definition}


\item 设$r > 0, a_1, a_2, \cdots, a_n$不全相等, 证明:
\[
M_q(r) > M_q\left(\frac{r}{2}\right).
\]


\begin{proof}

应用C-S不等式即可:
\[
M_q\left(\frac{r}{2}\right) = \left(\sum q_i a_i^{\frac{r}{2}}\right)^{\frac{2}{r}} = \left(\sum \sqrt{q_i} \sqrt{q_i} a_i^{\frac{r}{2}}\right)^{\frac{2}{r}} < \left(\sum q_i \sum q_i a_i^{r}\right)^{\frac{1}{r}} = M_q(r).
\]

\end{proof}


\item 证明:
\[
G_q = \lim_{r \to 0+} M_q(r).
\]


\begin{proof}

当$a_i >0 \ (1 \leq i \leq n)$时, 有
\[
a_i^r = e^{r \ln a_i} = 1 + r \ln a_i + o(r).
\]
故
\[
M_q(r) =  \left( \sum q_i a_i^r \right)^{\frac{1}{r}} = e^{\frac{1}{r} \ln \sum q_i a_i^r} = exp\left[ \frac{1}{r} \ln \sum q_i (1 + r \ln a_i + o(r)) \right]
\]
\[
= exp\left[ \frac{1}{r} \ln (1 + r \sum q_i \ln a_i + o(r)) \right].
\]
那么
\[
\lim_{r \to 0+}M_q(r) = exp\left( \sum q_i \ln a_i \right) = G_q.
\]

\end{proof}


\item 设$a_1, a_2, \cdots, a_n$不全相等, 则$G_q < M_q(1)$. 即:
\[
a_1^{q_1} a_2^{q_2} \cdots a_n^{q_n} < q_1 a_1 + \cdots + q_n a_n.
\]
等号成立的条件为$a_1 = a_2 = \cdots = a_n$.


\begin{proof}

根据以上二问:
\[
M_q(1) > M_q(\frac{1}{2}) >  M_q(\frac{1}{2^2}) > \cdots > M_q(\frac{1}{2^k}) > \lim_{r \to 0+} M_q(r) = G_q.
\]

\end{proof}



\item 已知椭圆$\frac{x^2}{a^2} + \frac{y^2}{b^2} =1 (a, b >0)$的周长:
\[
s = 4 \int_0^{\frac{\pi}{2}} \sqrt{a^2 \sin^2 t + b^2 \cos^2 t} dt,
\]
证明: $\pi(a+b) \leq s \leq \pi \sqrt{2a^2 + 2b^2}$. Tip: 利用Cauchy-Schwarz不等式积分/离散形式.


\begin{proof}


利用积分形式的Schwarz不等式得到上界估计:
\[
s = 4 \int_0^{\frac{\pi}{2}} \sqrt{a^2 \sin^2 t + b^2 \cos^2 t} dt \leq 4 \left( \int_0^{\frac{\pi}{2}} (a^2 \sin^2 t + b^2 \cos^2 t) dt  \right)^{\frac{1}{2}}  \left( \int_0^{\frac{\pi}{2}} dt \right)^{\frac{1}{2}} 
\]
\[
= 4 \left[ \frac{\pi}{4} (a^2 + b^2) \right]^{\frac{1}{2}} \left( \frac{\pi}{2} \right)^{\frac{1}{2}} = \pi \sqrt{2a^2 + 2b^2}.
\]


左边的不等式则利用离散形式的C-S不等式:
\[
\sqrt{a^2 \sin^2 t + b^2 \cos^2 t} = \sqrt{a^2 \sin^2 t + b^2 \cos^2 t} \cdot \sqrt{\sin^2 t + \cos^2 t} 
\]
\[
\geq a \sin t \cdot \sin t + b \cos t \cdot \cos t = a \sin^2 t + b \cos^2 t,
\]
对两边积分再乘4得到积分下界:
\[
s \geq 4 \int_0^{\frac{\pi}{2}} (a \sin^2 t + b \cos^2 t)dt = \pi (a + b).
\]


\end{proof}



\item 设$f(x)$在区间$[a, b]$上可积. 对区间$[a, b]$作等距分割: 
\[
\pi: a = x_0 < x_1 < \cdots < x_n = b, \ \Delta x_i = \frac{b-a}{n}. 
\]
推导此情形下的抛物线法公式.


{\bf 抛物线公式}
\[
\int_a^b f(x) dx \approx S_n = \sum_{i=1}^n \int_{x_{i-1}}^{x_i} (\alpha_i x^2 + \beta_i x + \gamma_i) dx,
\]
这里选择系数$\alpha_i, \beta_i, \gamma_i \ (i = 1, 2, \cdots, n)$, 使得抛物线$y = \alpha_i x^2 + \beta_i x + \gamma_i$经过以下三点
\[
(x_{i-1}, f(x_{i-1})), \ \left( \frac{x_i + x_{i-1}}{2}, f \left(  \frac{x_i + x_{i-1}}{2} \right) \right), \ (x_i, f(x_i)).
\]
计算得到
\begin{align*}
\int_{x_{i-1}}^{x_i} (\alpha_i x^2 + \beta_i x + \gamma_i) dx  = & \frac{1}{3} \alpha_i (x_i^3 - x_{i-1}^3) + \frac{1}{2} \beta_i (x_i^2 - x_{i-1}^2) + \gamma_i (x_i - x_{i-1}) \\
 = & \left[ \frac{1}{3} \alpha_i (x_i^2 + x_i x_{i-1} + x_{i-1}^2) + \beta_i \frac{x_i + x_{i-1}}{2} + \gamma_i \right] \Delta x_i\\
 = & \frac{1}{6} \left[ 2 \alpha_i (x_i^2 + x_i x_{i-1} + x_{i-1}^2) + 3 \beta_i (x_i + x_{i-1}) + 6 \gamma_i \right] \Delta x_i\\
 = &  \frac{1}{6} \{ (\alpha_i x_i^2 + \beta_i x_i + \gamma_i) + 4 \left[ \alpha_i  \left(  \frac{x_i + x_{i-1}}{2} \right)^2 + \beta_i  \left(  \frac{x_i + x_{i-1}}{2} \right) + \gamma_i \right] \\
& + (\alpha_i x_{i-1}^2 + \beta_i x_{i-1} + \gamma_i) \} \Delta x_i \\
 = & \frac{1}{6} \left[ f(x_i) + 4 f \left(  \frac{x_i + x_{i-1}}{2} \right) + f(x_{i-1}) \right] \Delta x_i.
\end{align*}
因此计算得到
\[
S_n = \frac{b-a}{6n} \sum_{i=1}^n \left[ f(x_i) + 4 f \left(  \frac{x_i + x_{i-1}}{2} \right) + f(x_{i-1}) \right].
\]
可以看到
\[
S_n = \frac{2}{3} R_n + \frac{1}{3} T_n.
\]



\item 设$f \in C^2[0, h]$, 证明: 存在$\xi \in [0, h]$使得
\[
\int_0^h f(x) dx = \frac{h}{2} [f(0) + f(h)] - \frac{1}{12} f''(\xi) h^3. 
\]


\begin{proof}


作两次分部积分得:
\[
\int_0^h f(x) dx = \int_0^h f(x) d \left( x - \frac{h}{2} \right) = f(x) \left( x - \frac{h}{2} \right)\Big|_0^h - \int_0^h f'(x)  \left( x - \frac{h}{2} \right)dx
\]
\[
= \frac{h}{2} [f(0) + f(h)] - \frac{1}{2} \int_0^h f'(x) d (x^2 - hx)  = \frac{h}{2} [f(0) + f(h)] + \frac{1}{2} \int_0^h f''(x) (x^2 - hx) dx. \eqno{(1)}
\]
由于$x^2 - hx$不变号, 对右边的积分用第一中值定理得
\[
\frac{1}{2} \int_0^h f''(x) (x^2 - hx)  dx = \frac{1}{2} f''(\xi) \int_0^h (x^2 - hx)  dx = - \frac{1}{12} f''(\xi) h^3.
\]
也可以用微积分基本定理证明.
\end{proof}


\item 设$f \in C^4[0, h]$, 证明: 存在$\xi \in [0, h]$使得
\[
\int_0^h f(x) dx = \frac{h}{2} [f(0) + f(h)] - \frac{h^2}{12} [f'(h) - f'(0)] + \frac{1}{720} f^{(4)} (\xi) h^5.
\]

\begin{proof}

从上问中(1)式继续推导:
\begin{align*}
\int_0^h f(x) dx - \frac{h}{2} [f(0) + f(h)] & = \frac{1}{2} \int_0^h f''(x)(x^2 - hx) dx \\
& = - \frac{h^2}{12} \int_0^h f''(x) dx + \frac{1}{2} \int_0^h f''(x) \left(x^2 - hx + \frac{h^2}{6} \right) dx \\
& = - \frac{h^2}{12} [f'(h) - f'(0)] + \frac{1}{2} \int_0^h f''(x) d \left(\frac{x^3}{3} - \frac{hx^2}{2} + \frac{h^2x}{6} \right) \\
& =  - \frac{h^2}{12} [f'(h) - f'(0)] - \frac{1}{2} \int_0^h f'''(x) \left(\frac{x^3}{3} - \frac{hx^2}{2} + \frac{h^2x}{6} \right) dx \\
& = - \frac{h^2}{12} [f'(h) - f'(0)] - \frac{1}{2} \int_0^h f'''(x) d \left(\frac{x^4}{12} - \frac{hx^3}{6} + \frac{h^2x^2}{12} \right) \\
& = - \frac{h^2}{12} [f'(h) - f'(0)] + \frac{1}{24} \int_0^h f^{(4)} (x) [x^2 (x-h)^2] dx,
\end{align*}
在最后一个积分式中, 利用因子$x^2(x - h)^2$不变号, 利用积分第一中值定理即得所需等式.


\end{proof}


\item 设$f(x)$在$[0, 1]$上可导, 且当$x \in (0, 1)$时, $0 < f'(x) < 1, f(0) = 0$. 试证:
\[
\left( \int_0^1 f(x) dx \right)^2 > \int_0^1 f^3(x) dx.
\]

\begin{proof}

\noindent
{\bf 证一:} 令
\[
F(x) = \left( \int_0^x f(t) dt \right)^2 - \int_0^x f^3(t) dt.
\]
则$F(0) = 0$. 只需证在$(0, 1)$内$F'(x) > 0$即可. 事实上, 
\[
F'(x) = f(x) \left[ 2 \int_0^x f(t) dt - f^2(x) \right].
\]
已知$f(0) = 0, 0 < f'(x) < 1$, 故当$x \in (0, 1)$时$f(x) > 0$. 记
\[
g(x) = 2 \int_0^x f(t) dt - f^2(x),
\]
则$g(0) = 0$, 且
\[
g'(x) = 2 f(x) [1 - f'(x)] > 0 \Rightarrow g(x) > 0.
\]


\noindent
{\bf 证二:} 将问题转化为证明
\[
\frac{\left( \int_0^1 f(x) dx \right)^2}{ \int_0^1 f^3(x) dx } > 1
\]
令
\[
F(x) = \left( \int_0^x f(t) dt \right)^2 \ \ G(x) =  \int_0^x f^3(t) dt. 
\]
利用Cauchy中值定理
\[
\frac{\left( \int_0^1 f(x) dx \right)^2}{ \int_0^1 f^3(x) dx }  = \frac{F(1) - F(0)}{G(1) - G(0)} = \frac{F'(\xi)}{G'(\xi)} = \frac{2f(\xi) \int_0^{\xi} f(t) dt}{f^3(\xi)} = \frac{2 \int_0^{\xi} f(t) dt}{f^2(\xi)} \ (0 < \xi < 1)
\]
\[
=  \frac{2 \int_0^{\xi} f(t) dt - 2 \int_0^0 f(t)dt}{f^2(\xi) - f^2(0)} = \frac{2f(\eta)}{2f(\eta)f'(\eta)} = \frac{1}{f'(\eta)} > 1 \ (0 < \eta < \xi < 1).
\]

\end{proof}


\end{enumerate}

\clearpage

\section{5月13--19日 · Week 12}

\subsection*{题目}

\begin{enumerate}

\item 设$x \geq m \in \mathbb N^*$时, $f$为非负递减函数, 则极限
\[
\lim_{n \to \infty} \left( \sum_{k=m}^n f(k) - \int_m^n f(x) dx \right) = \alpha \eqno{(4)}
\]
存在, 且$0 \leq \alpha \leq f(m)$. 进一步地, 如果$\lim\limits_{x \to +\infty} f(x) =0$, 那么
\[
\Big| \sum_{k=m}^{[\xi]} f(k) - \int_m^{\xi} f(x) dx - \alpha \Big| \leq f(\xi -1), \eqno{(5)}
\]
这里$\xi \geq m+1$.

\item 设正项数列$\{ a_n \}$单调递减, 则$\lim_{n \to \infty} a_n = 0$的充分必要条件是正项级数$\sum_{n=1}^{\infty} \left( 1 - \frac{a_{n+1}}{a_n} \right)$发散.

\item 设$a_n >0$, $ S_n = a_1 + \cdots + a_n$, 证明:
\[
\sum_{n=1}^{\infty} \frac{a_n}{S_n^2} < +\infty.
\]

\item 设$\sum_{n=1}^{\infty} a_n$为正项级数, $p>1$, 证明: 无论$\sum_{n=1}^{\infty} a_n$收敛与否, 级数$\sum_{n=1}^{\infty} \frac{a_n}{S_n^p}$总收敛.

\item 设$p>0, q>0$. 当$p, q$取何值时, 级数
\[
\sum_{n=1}^{\infty} \frac{p(p+1)\cdots (p+n-1)}{n!} \frac{1}{n^q}
\]
收敛. (用Raabe判别法)

\item 设$\sum_{n=1}^{\infty} a_n$是递减的正项发散级数, 证明当$n \to \infty$时
\[
\frac{a_1 + a_3 + \cdots + a_{2n-1}}{a_2 + a_4 + \cdots + a_{2n}} \to 1.
\]

\item 设
\[
a_n = \left( 1 - \frac{p \ln n}{n} \right)^n,
\]
证明级数$\sum_{n=1}^{\infty} a_n$收敛.

\end{enumerate}

\clearpage

\subsection*{原稿答案}

\noindent\emph{以下为按题干归入本周的原稿部分参考答案。}

\begin{enumerate}


\item 设正项数列$\{ a_n \}$单调递减, 则$\lim_{n \to \infty} a_n = 0$的充分必要条件是正项级数$\sum_{n=1}^{\infty} \left( 1 - \frac{a_{n+1}}{a_n} \right)$发散.


\begin{proof}

记$b_n = 1 -  \frac{a_{n+1}}{a_n}, n \geq 1$. 由于正数列$\{ a_n \}$单调递减, 因此$\lim_{n \to \infty} a_n = a$存在. 


若$a >0$, 则$b_n \leq \frac{a_n - a_{n+1}}{a}$. 由于$\sum_{n=1}^{\infty} (a_n - a_{n+1}) = \lim_{n \to \infty} (a_1 - a_n) = a_1 - a$, 则级数收敛(逆否命题).


若$a = 0$, 则
\[
\sum_{k = n+1}^{n+p} b_k = \sum_{k = n+1}^{n+p} \frac{a_k - a_{k+1}}{a_k} \geq \sum_{k = n+1}^{n+p}  \frac{a_k - a_{k+1}}{a_{n+1}} = 1 - \frac{a_{n+p+1}}{a_{n+1}}.
\]
当固定$n$, 并取$p$充分大时, 上式最后一项$> \frac{1}{2}$. 利用Cauchy收敛法则说明级数发散.

\end{proof}


\item 设$a_n >0$, $ S_n = a_1 + \cdots + a_n$, 证明:
\[
\sum_{n=1}^{\infty} \frac{a_n}{S_n^2} < +\infty.
\]

\begin{proof}


由于$\sum_{n=1}^{\infty}  \frac{a_n}{S_n^2}$是正项级数, 只需证明其部分和有界. 事实上, 对任意的正整数$N$,  有
\[
\sum_{n=2}^{N} \frac{a_n}{S_n^2} = \sum_{n=2}^{N} \frac{S_n - S_{n-1}}{S_n^2} \leq \sum_{n=2}^{N} \frac{S_n - S_{n-1}}{S_nS_{n-1}}
\]
\[
= \sum_{n=2}^{N} \left( \frac{1}{S_{n-1}} - \frac{1}{S_n} \right) = \frac{1}{S_1} - \frac{1}{S_N} < \frac{1}{a_1}.
\]
\end{proof}


\item 设$\sum_{n=1}^{\infty} a_n$为正项级数, $p>1$, 证明: 无论$\sum_{n=1}^{\infty} a_n$收敛与否, 级数$\sum_{n=1}^{\infty} \frac{a_n}{S_n^p}$总收敛.

\begin{proof}

利用Lagrange中值定理, 设$0 < a < b$, 在$[a, b]$上考虑函数$x^{1-p}$, 则存在$\xi \in (a, b)$使得
\[
\frac{1}{a^{p-1}} - \frac{1}{b^{p-1}} = \frac{p-1}{\xi^p} \cdot (b-a) > ( p - 1) \cdot \frac{b-a}{b^p}.
\]
取$a = S_{n-1}, b = S_n$, 得到
\[
\frac{a_n}{S_n^p} = \frac{S_n - S_{n-1}}{S_n^p} < \frac{1}{p-1} \left( \frac{1}{S_{n-1}^{p-1}} - \frac{1}{S_n^{p-1}} \right).
\]
因此
\[
\sum_{k=1}^{n} \frac{a_k}{S_k^p} = \frac{1}{a_1^{p-1}} + \sum_{k=2}^{n} \frac{a_k}{S_k^p} < \frac{1}{a_1^{p-1}} + \frac{1}{p-1} \sum_{k=2}^{n} \left( \frac{1}{S_{k-1}^{p-1}} - \frac{1}{S_k^{p-1}} \right)
\]
\[
< \frac{1}{a_1^{p-1}} + \frac{1}{p-1} \cdot \frac{1}{a_1^{p-1}} = \frac{p}{p-1} \cdot \frac{1}{a_1^{p-1}}.
\]
因此级数收敛.

\end{proof}


\item 设$p>0, q>0$. 当$p, q$取何值时, 级数
\[
\sum_{n=1}^{\infty} \frac{p(p+1)\cdots (p+n-1)}{n!} \frac{1}{n^q}
\]
收敛. (用Raabe判别法)


\begin{proof}

利用Raabe判别法,
\[
\frac{a_n}{a_{n+1}} = \frac{n+1}{n+p} \left( 1 + \frac{1}{n}  \right)^q =  \left( 1 + \frac{1-p}{n+p}  \right)  \left( 1 + \frac{q}{n} + o \left( \frac{1}{n} \right) \right)
\]
\[
= 1 + \frac{q}{n} + \frac{1-p}{n+p} + o \left( \frac{1}{n} \right)= 1 + \frac{q+1-p}{n} + o \left( \frac{1}{n} \right).
\]
故原级数当$q > p$时收敛, 当$q < p$时发散, 当$q =p$时无法判断.

\end{proof}


\item 设$\sum_{n=1}^{\infty} a_n$是递减的正项发散级数, 证明当$n \to \infty$时
\[
\frac{a_1 + a_3 + \cdots + a_{2n-1}}{a_2 + a_4 + \cdots + a_{2n}} \to 1.
\]

\begin{proof}

记分子分母为$A, B$. 显然$A \geq B$且
\[
A - a_1 = a_3 + a_5 + \cdots + a_{2n-1} \leq B.
\]
又
\[
2B \geq B + A - a_1 = \sum_{k=2}^n a_k \to \infty \ (n \to \infty).
\]
故
\[
\Big| \frac{A}{B} -1 \Big| = \frac{|A-B|}{B} \leq \frac{a_1}{B} \to 0.
\]

\end{proof}


\item 设
\[
a_n = \left( 1 - \frac{p \ln n}{n} \right)^n,
\]
证明级数$\sum_{n=1}^{\infty} a_n$收敛.

\begin{proof}

由于
\[
\lim_{n \to \infty} \ln (n^p a_n) = \lim_{n \to \infty} \ln \left[  n^p  \left( 1 - \frac{p \ln n}{n} \right)^n \right]
\]
\[
= \lim_{n \to \infty} \left[ p \ln n + n \ln \left( 1 - \frac{p \ln n}{n} \right) \right]
\]
\[
= \lim_{n \to \infty} \frac{1}{\frac{1}{n}} \left[ - \frac{1}{n} p \ln \frac{1}{n} + \ln \left( 1 + \frac{p}{n} \ln \frac{1}{n} \right) \right]
\]
\[
= \lim_{x \to 0} \frac{1}{x} [ \ln(1 + px \ln x) - px \ln x] =0.
\]
故$a_n \sim 1/n^p$, 故当$p > 1$时收敛, 当$ p \leq 1$时发散.

\end{proof}

\end{enumerate}

\clearpage

\section{5月27日--6月2日 · Week 14}

\subsection*{题目}

\begin{enumerate}

\item 设$\sum_{n=1}^{\infty} u_n(x)$为$[a, b]$上可导的函数项级数, 且$\sum_{n=1}^{\infty} u_n'(x)$的部分和函数列在$[a, b]$上一致有界. 证明: 若$\sum_{n=1}^{\infty} u_n(x)$在$[a, b]$上逐点收敛, 则一致收敛.

\item 设函数列$\{f_n(x) \} (n \geq 0)$在区间$I$上有定义, 且$\forall x \in I$满足

(1) $|f_0(x)| \leq M$;

(2) $\sum_{n=0}^m |f_n(x) - f_{n+1}(x)| \leq M, m = 0, 1, 2, \cdots$, 其中$M$为常数. 证明: 若级数$\sum_{n=0}^{\infty} b_n$收敛, 则级数$\sum_{n=0}^{\infty} b_n f_n(x)$在$I$上一致收敛.

Tip: 使用Abel变换和Cauchy收敛准则.

\item 设$\{ u_n(x) \}$是$[a, b]$上正的递减函数列$(u_{n+1}(x) \leq u_n(x))$, 且处处收敛于0. 每个$u_n(x)$都是$[a, b]$上的递增函数. 证明: 级数$\sum_{n=1}^{\infty} (-1)^{n-1} u_n(x)$在$[a, b]$上一致收敛.

\item 证明函数列$\left( 1 + \frac{x}{n} \right)^n \ (n \geq 1)$一致收敛.

\item 给定函数列
\[
f_n(x) = \frac{x (\ln n)^{\alpha}}{n^x}, \ \ n \geq 2.
\]
问当$\alpha$取何值时, $\{ f_n(x) \}$在$[0, +\infty)$上一致收敛.

\item 设对每个$n$, 函数$u_n(x)$在$x = c$处左连续, 又设$\sum_{n=1}^{\infty} u_n(c)$发散. 那么对任意$\delta >0$, $\sum_{n=1}^{\infty} u_n(x)$在区间$(c-\delta, c)$上必不一致收敛.

\end{enumerate}

\clearpage

\subsection*{原稿答案}

\noindent\emph{原稿未附这一周对应的参考答案。}

\clearpage

\section{6月3--9日 · Week 15}

\subsection*{题目}

\begin{enumerate}

\item 设$f_n(x)$在$(-\infty, +\infty)$上一致连续, 且$f_n(x)$在$(-\infty, +\infty)$上一致收敛于$f(x)$. 证明$f(x)$在$(-\infty, +\infty)$上一致连续.

\item 设$\{ f_n(x) \}$是$[a, b]$上的连续函数列, 且$f_n(x)$一致收敛于$f(x)$. 若$\{ x_n \} \subset [a, b], x_n \to x_0 \ (n \to \infty)$, 证明:
\[
\lim_{n \to \infty} f_n(x_n) = f(x_0).
\]

\item 证明级数$\sum_{n=-\infty}^{\infty} \frac{1}{(n-x)^2}$当$x \notin \mathbb Z$时收敛, 周期为1, 并且当$x \neq \mathbb Z$时和函数连续.

\item 若对$R>0$, 成立
\[
\Big| \sum_{k=0}^n a_k R^k \Big| \leq M \ \ (n=0, 1, 2, \cdots)
\]
则
\[
(1) \ \  \sum_{n=0}^{\infty} a_n x^n \ (0 < x < R)\mbox{收敛}; \ \ \ \ (2) \ \  \Big| \sum_{n=0}^{\infty} a_n x^n \Big| \leq M \ (0 < x < R).
\]

\item 求下述级数的和:
\[
\sum_{n=1}^{\infty} \frac{(-1)^{n-1}}{(2n-1)(2n+1)}.
\]

\item 求和
\[
1 - \frac{1}{2} + \frac{1}{4} - \frac{1}{5} + \frac{1}{7} -  \cdots.
\]

\end{enumerate}

\clearpage

\subsection*{原稿答案}

\paragraph{第 1 题}
\begin{proof}

根据一致收敛条件, $\forall \epsilon >0, \exists N \in \mathbb N^*$, 使得当$n >N$时, 对$\forall x \in (-\infty, +\infty)$, 成立
\[
|f(x) - f_n(x)| < \frac{\epsilon}{3}.
\]
不妨取$n = N+1$(可以固定$n$来考虑). 对任意$x_1, x_2 \in (-\infty, +\infty)$, 有
\[
|f(x_1) - f(x_2)| \leq |f(x_1) - f_n(x_1)| + |f_n(x_1) - f_n(x_2)| + |f(x_2) - f_n(x_2)| \leq \frac{2\epsilon}{3} + |f_n(x_1) - f_n(x_2)|.
\]
由于$f_n(x)$一致连续($n$固定), 因此对上述$\epsilon >0, \exists \delta >0$, 使得当$|x_1 - x_2|<\delta$时, 
\[
|f_n(x_1) - f_n(x_2)| < \frac{\epsilon}{3}.
\]

\end{proof}

\paragraph{第 2 题}
\begin{proof}

根据一致收敛性, $\forall \epsilon >0, \exists N_1$(与$x$无关), 当$n >N_1$时, 对$\forall x \in [a, b]$, 有
\[
|f_n(x) - f(x)| < \frac{\epsilon}{2}.
\]
特别地
\[
|f_n(x_n) - f(x_n)| < \frac{\epsilon}{2}.
\]
根据连续性定理, $f(x)$在$[a, b]$上连续, 即$\exists \delta > 0$, 使得当$|x - x_0| < \delta$时,
\[
|f(x) - f(x_0)| < \frac{\epsilon}{2}.
\]
所以当$|x - x_0| < \delta$时,
\[
|f_n(x_n) - f(x_0)| \leq |f_n(x_n) - f(x_n)| + |f(x_n) - f(x_0)| \leq \frac{\epsilon}{2} + \frac{\epsilon}{2} = \epsilon.
\]
即证.

\end{proof}

\paragraph{第 3 题}
\begin{proof}

由于
\[
\sum_{n=-\infty}^{\infty} \frac{1}{(n-x)^2} = \sum_{n=0}^{\infty} \frac{1}{(n-x)^2} + \sum_{n=1}^{\infty} \frac{1}{(-n-x)^2}.
\]
当$x \notin \mathbb Z$, $n \to \infty$时
\[
\frac{1}{(n-x)^2} \sim \frac{1}{n^2}, \ \  \frac{1}{(-n-x)^2} \sim \frac{1}{n^2}.
\]
故级数收敛. 又有
\[
f(x+1) = \sum_{n=-\infty}^{\infty} \frac{1}{(n-1-x)^2}  = \sum_{k=-\infty}^{\infty} \frac{1}{(k-x)^2} \equiv f(x). 
\]
其连续性只须在$(0, 1)$内证明. 由于在$(0, 1)$内
\[
\Big| \frac{1}{(n-x)^2} \Big| \leq \frac{1}{(n-1)^2},  \ \ \Big| \frac{1}{(-n-x)^2} \Big| = \frac{1}{(n+x)^2} \leq \frac{1}{n^2}.
\]
故级数在$(0, 1)$一致收敛, 因此和函数连续.

\end{proof}

\paragraph{第 4 题}
\begin{proof}

(1) 当$0 < x < R$时, 数列$\{ \left( \frac{x}{R} \right)^n \}$递减趋于0, 再注意到级数$\sum_{n=0}^{\infty} a_n R^n$的部分和数列有界, 因此
\[
\sum_{n=0}^{\infty} a_n x^n = \sum_{n=0}^{\infty} a_n R^n \cdot \left( \frac{x}{n} \right)^n,
\]
利用数项级数的Abel判别法知级数处处收敛.

(2) 记$S_n = \sum_{k=0}^n a_k R^k$, 则利用Abel变换
\[
\sum_{k=0}^n a_k x^k = \sum_{k=0}^n a_k R^k \left( \frac{x}{R} \right)^k = \sum_{k=0}^{n-1} S_k \left( \left( \frac{x}{R} \right)^k - \left( \frac{x}{R} \right)^{k+1} \right) + S_n \left( \frac{x}{R} \right)^n
\]
\[
\leq  M \left[ 1 -  \left( \frac{x}{R} \right)^n \right] + S_n \left( \frac{x}{R} \right)^n.
\]
两边取绝对值, 得
\[
\Big| \sum_{k=0}^n a_k x^k \Big| \leq M \left[ 1 + 2 \Big| \frac{x}{R} \Big|^n \right].
\]
令$n \to \infty$即证.

\end{proof}

\paragraph{第 5 题}
\begin{proof}
\

\noindent
{\bf 解一:} 令
\[
f(x) = \sum_{n=1}^{\infty} \frac{(-1)^{n-1}}{(2n-1)(2n+1)} x^{2n+1}.
\]
易得收敛半径$R =1$, 且$f(0) = 0$. 逐项求导得到:
\[
f'(x) = \sum_{n=1}^{\infty} \frac{(-1)^{n-1}}{2n-1} x^{2n} \ (|x| < 1).
\]
再令
\[
g(x) =  \sum_{n=1}^{\infty} \frac{(-1)^{n-1}}{2n-1} x^{2n-1} \ (|x| < 1).
\]
则$g(0) = 0$. 因此对$|x| < 1$,
\[
g(x) = \int_0^x g'(t) dt =  \int_0^x \sum_{n=1}^{\infty} (-1)^{n-1} x^{2n-2} dx = \int_0^x \sum_{n=0}^{\infty} (-1)^{n} x^{2n} dx = \int_0^x \frac{1}{1+t^2} dt = \arctan x.
\]
因此$f'(x) = x \arctan x$, 于是
\[
f(x) = \int_0^x f'(t) dt = \int_0^x t \arctan t dt = \frac{x^2}{2} \arctan x - \frac{x}{2} + \frac{1}{2} \arctan x \ (|x| < 1).
\]
验证级数在$x=1$处收敛, 再利用Abel引理,
\[
\sum_{n=1}^{\infty} \frac{(-1)^{n-1}}{(2n-1)(2n+1)} = \lim_{x \to 1-} f(x) = \frac{\pi -2}{4}.
\]


\noindent
{\bf 解二:} 作分解:
\[
\sum_{n=1}^{\infty} \frac{(-1)^{n-1}}{(2n-1)(2n+1)} = \frac{1}{2} \sum_{n=1}^{\infty} (-1)^{n-1} \left[ \frac{1}{2n-1} - \frac{1}{2n+1} \right]
\]
\[
= \frac{1}{2} \sum_{n=1}^{\infty} \frac{(-1)^{n-1}}{2n-1} - \frac{1}{2} \sum_{n=1}^{\infty} \frac{(-1)^{n-1}}{2n+1}
\]
\[
= \frac{1}{2} \sum_{n=1}^{\infty} \frac{(-1)^{n-1}}{2n-1} + \frac{1}{2} \sum_{k=2}^{\infty} \frac{(-1)^{k-1}}{2k-1}.
\]
利用逐项积分, 对任意$x \in (-1, 1)$, 有
\[
\arctan x = \int_0^x \frac{dt}{1+t^2} = \int_0^x \sum_{k=1}^{\infty} (-1)^{k-1} t^{2k-2} dt = \sum_{k=1}^{\infty} (-1)^{k-1} \frac{x^{2k-1}}{2k-1}.
\] 
再根据Abel引理,
\[
\sum_{k=1}^{\infty} \frac{(-1)^{k-1}}{2k-1} = \lim_{x \to 1-} \arctan x = \frac{\pi}{4}.
\]
综上可得数项级数的和为$ \frac{\pi -2}{4}$.

\end{proof}

\paragraph{第 6 题}
\begin{proof}

由
\[
\frac{1}{1 + x + x^2} = (1-x)(1 + x^3 + x^6 + \cdots ) =  1 -x + x^3 - x^4 + \cdots.
\]
考虑
\[
f(x) = \int_0^x \frac{dt}{1 + t + t^2} = x - \frac{x^2}{2} + \frac{x^4}{4} - \frac{x^5}{5} + \cdots, 
\]
则该级数和为$f(1)$. 由于
\[
\int_0^x \frac{dt}{1 + t + t^2} = \frac{2}{\sqrt 3} \arctan \left( \frac{x + 1/2}{\sqrt 3 /2} \right) - \frac{2}{3} \arctan \left( \frac{1}{\sqrt 3} \right), \ \ 0 \leq x < 1.
\]
由Abel定理知
\[
f(1) = \lim_{x \to 1-} f(x) = \frac{2}{\sqrt 3} \arctan \left( \frac{3/2}{\sqrt 3 /2} \right) - \frac{2}{3} \arctan \left( \frac{1}{\sqrt 3} \right) = \frac{\pi}{3 \sqrt 3}.
\]
\end{proof}

\end{document}

